% !TeX program = lualatex
% 共通の組版設定。本文と証明は各論文の節ファイルに一度だけ置く。
\documentclass[11pt,a4paper]{ltjsarticle}
\usepackage[margin=24mm]{geometry}
\usepackage{amsmath,amssymb,amsthm,mathtools,booktabs,longtable}
\usepackage{mathrsfs}
\usepackage[haranoaji]{luatexja-preset}
\usepackage[unicode,hidelinks]{hyperref}
\setlength{\emergencystretch}{3em}
\hypersetup{bookmarksnumbered=true}
\pdfstringdefDisableCommands{\def\qgr{qgr}\def\HH{HH}\def\Perf{Perf}}
\providecommand{\PublicationVersion}{publication-v32}
\newtheorem{theorem}{定理}[section]
\newtheorem{proposition}[theorem]{命題}
\newtheorem{lemma}[theorem]{補題}
\newtheorem{corollary}[theorem]{系}
\theoremstyle{definition}\newtheorem{definition}[theorem]{定義}
\newtheorem{assumption}[theorem]{入力・仮定}
\theoremstyle{remark}\newtheorem{remark}[theorem]{注意}
\renewcommand{\proofname}{証明}
\newcommand{\HH}{\operatorname{HH}}
\newcommand{\qgr}{\operatorname{qgr}}
\newcommand{\coh}{\operatorname{coh}}
\newcommand{\Perf}{\operatorname{Perf}}
\newcommand{\Ext}{\operatorname{Ext}}
\newcommand{\Hom}{\operatorname{Hom}}
\newcommand{\Spec}{\operatorname{Spec}}
\newcommand{\Spf}{\operatorname{Spf}}
\newcommand{\GL}{\operatorname{GL}}
\newcommand{\gldim}{\operatorname{gldim}}
\newcommand{\OO}{\mathcal O}
\renewcommand{\AA}{\mathbb A}
\newcommand{\PP}{\mathbb P}

% 各原稿の引用だけを共通文献データから組版する。
\makeatletter
\newcommand{\UsePublicationReferences}[1]{%
  \@for\pub@reference:=#1\do{%
    \expandafter\def\csname pubref@\pub@reference\endcsname{}%
  }%
}
\makeatother

\title{三乗根量子三次曲面の正常形と形式変形}
\author{}
\date{2026年11月18日\quad publication-v32}
\begin{document}
\maketitle
\begin{abstract}
原始三乗根で定義される量子 $\mathfrak{so}_3$ 包絡と、その正則中心三次商を研究する。
中央境界を固定したPBW拡大の対角非零部分を七次元の正常形で記述し、
境界の三捩れ平行移動を動かす八変数の形式表示族を構成する。
滑らかな三次中心と反標準楕円分岐因子に沿う指定局所order模型の下で、
Hochschild変形空間の次元は8、三次障害空間は零となる。
Artin底上の反標準環の再構成、二次包絡の平坦性の障害、
形式表示族と適合したAbel変形の比較を分けて証明する。
全H3モジュライの代数化・開性・降下を仮定なしに結論しない。
反標準環からの二次包絡の一意性を用い、AOUの無限位数のblow-up構成との
限定された非同一性も示す。
\end{abstract}

\section{序論}\label{sec:e6-introduction}
非可換三次曲面を中心超曲面として与える表示と、構造対象を持つ圏から
反標準環を再構成する記述を、具体的な変形問題の上で比較する。
基礎体 $k$ は標数0の代数閉体、$q$ は原始三乗根とする。
基本的な包絡は
\[
 B_{\mathbf a}=k\langle t,x,y,z\rangle/
 (yz-qzy-tx-at^2,\;zx-qxz-ty-bt^2,\;xy-qyx-tz-ct^2,\;[t,x],[t,y],[t,z])
\]
であり、中心三次元の直線 $kG$ を選んで $R=B_{\mathbf a}/(G)$ を作る。
この表示のパラメータ数、圏のHochschild接空間、全モジュライの局所次元は、
比較定理を介して初めて結び付く。

\subsection{中心結果と証明の依存関係}
第一の結果は、固定位数三・中央境界・対角非零のPBW部分の正常形である。
対称一次補正の非対角成分を平行移動で消し、対角成分を有限\'etale底変更で規格化すると、
上の $B_{\mathbf a}$ に帰着する（定理\ref{thm:normal}）。
残余群は偶符号変更と巡回置換からなる $\mathfrak A_4$ であり、
中心三次直線の四次元と $(a,b,c)$ の三次元を合わせて七次元を得る。
ここでの中心・PBW・同型計算は、H3認識の仮定とは独立した代数的主張である。

第二の結果は、次の幾何学的仮定の下での変形空間の計算である。
\begin{theorem}[Hochschild変形空間：定理\ref{thm:HH8}]
滑らかな射影三次曲面 $C$ と滑らかな $D\in|-K_C|$ 上のorder $\Lambda$ が、
$C\setminus D$ ではAzumayaであり、$D$ に沿って\'etale局所的な完備化で
$k[[s,w]]\rtimes\mu_3$、$\zeta(s,w)=(s,\zeta w)$ とMorita同値であるとする。
このとき $\mathcal C=\Perf(C,\Lambda)$ に対し
\[
 \HH^0(\mathcal C)=k,\qquad \HH^1(\mathcal C)=0,\qquad
 \dim\HH^2(\mathcal C)=8,\qquad \HH^n(\mathcal C)=0\quad(n\ge3).
\]
通常のHochschild複体が制御する形式変形は無障害であり、
miniversal底は八次元形式円板となる。
\end{theorem}
証明は局所固定部分の計算と $T_C(-\log D)$ のコホモロジーによる。
具体的な中心商がこの局所order模型を満たすという認識は入力\ref{ass:order}で明示する。
接空間行列の階数計算から、その幾何学的認識を逆に推論しない。

第三の結果は、七次元部分を横断する実際の形式族の構成である。
中央三次の低次補正の一意性（補題\ref{lem:central-lift}）を使って
中心が跳ぶ三捩れ点でも指定切断を持ち上げ、八変数の同時形式表示を得る
（定理\ref{thm:eight-formal-central-lift}）。
適合したAbel変形との比較は第\ref{sec:artin-reconstruction}節と
第\ref{sec:artin-enhancement}節、代数的な表示stackは第\ref{sec:presentation-stack}節で扱う。
これらを全intrinsic H3の成分定理へ進めるために必要な追加条件は、
最後の定理\ref{thm:conditional-component}にまとめる。

\begin{center}
\begin{tabular}{p{.28\textwidth}p{.62\textwidth}}
\toprule
主張 & 適用範囲と証明箇所\\\midrule
七次元正常形 & 中央PBW・対角非零。定理\ref{thm:normal}。\\
八次元Hochschild変形 & 指定した滑らかな中心・分岐order。定理\ref{thm:HH8}。\\
全Artin形式表示 & 指定点、表示chartとenhancementの比較。第\ref{sec:artin-enhancement}節。\\
大域成分の一意性 & 代数的モジュライと相対比較の追加仮定付き。定理\ref{thm:conditional-component}。\\
AOUとの比較 & 無限位数のblow-up構成とのpointed反標準比較。第\ref{sec:as-surface-scope}節。\\
\bottomrule
\end{tabular}
\end{center}

\subsection{先行研究と本稿の位置づけ}
Etingof--Ginzburg \cite[Theorem 3.3.2]{EG} は、指定された次数の斉次CYポテンシャルへの
低次項の付加がPBW型CY変形を与え、非定数中心元が存在することを示した。
本稿はこの存在定理を用いる。中心が跳ぶ特殊点で、選んだ中心元を
正則に持ち上げる手順、表示の規格化、および圏の変形との比較を別に証明する。
楕円 $E_6$ という呼称は同文献の三次ポテンシャルの型に由来し、
本稿だけでWeyl群商による全モジュライ分類を得たという意味ではない。

Chekhov--Mazzocco--Rubtsov \cite[Definition 1.1, (1.5), Theorem 1.5]{CMR} は
普遍Painlev\'e代数とその中心特殊化を扱う。
同文献のDefinition 1.1には $q$ が根でないという仮定があるため、
その定理を三乗根へ無条件に適用しない。
表示式の具体的な特殊化と、本稿で直接証明する三乗根での中心増大を区別する。
一般化Sklyanin--Painlev\'e族の全同型類や結節因子との完全な既出性比較は未完であり、
本稿の各計算を網羅的な新規性確認済みとはしない。

局所Hochschild層はNegron--Schedler \cite{NS} の商の公式と整合し、
orderの大域的な捩れはFaber--Ingalls--Okawa--Satriano \cite{FIOS} の枠組みと比較する。
Brauer捩れを消して通常のroot stackと同一視する必要はない。
Abel変形と有向 $\mathbb Z$-代数の比較にはDe Deken--Lowen \cite{DeDekenLowen}、
Hochschild障害理論にはLowen--Van den Bergh \cite{LVdB,LowenObject} を用いる。
Rains \cite{Rains} の次数三の全パラメータ被覆は証明済み入力として使わない。
AS surfaceという語はAOU/Okawa--Uedaの定義と区別して扱い、
反標準二次包絡から得られる限定的な比較を第\ref{sec:as-surface-scope}節に置く。

\subsection{規約と研究の分担}
本稿のH3は、構造対象 $\OO$ を持つExt有限noetherian Abel圏で、
$\tau=S^{-1}[2]$ のheart保存、Gabriel--Krull rank-one torsion-free exceptionality、
Artin--Zhang豊富性、三次曲面の反標準Hilbert値を満たし、
反標準環が一次生成で、その二次包絡がquadratic Koszul AS4型14641である強いクラスを指す。
同型は構造対象を保つexact同値とする。必要な認識入力は\ref{ass:recognition}に列挙する。
この定義を一般の導来同値で割ったモジュライ問題へ読み替えない。

Fricke中心、Poisson臨界被覆、$D_4$周期とtrialityの証明は関連論文\cite{CompanionFricke}、
Vieta反復の境界特異点とPrymの証明は\cite{CompanionBoundary}に収める。
これらの主結果は、本稿の未完成の大域H3分類を前提としない。

\tableofcontents
\section{三乗根包絡と中心三次切断}\label{sec:seven}
$q^2+q+1=0$ を固定する。$\mathbf a=(a,b,c)$ に対し
\begin{equation}\label{eq:B}
B_{\mathbf a}=\frac{k\langle t,x,y,z\rangle}
{(yz-qzy-tx-at^2,\ zx-qxz-ty-bt^2,\ xy-qyx-tz-ct^2,\ [t,x],[t,y],[t,z])}
\end{equation}
と置く。$d=(q-1)^{-1}=-(q+2)/3$ とし、
\begin{align}
X_h&=x^3+q^2t^2x,\quad Y_h=y^3+q^2t^2y,\quad Z_h=z^3+q^2t^2z,\\
\Omega_{\mathbf a,h}&=xyz+\frac{q-1}{3}t(x^2+z^2)+\frac{q+2}{3}ty^2
                 +dt^2(ax+qby+cz).\label{eq:Omega}
\end{align}

\begin{remark}[普遍Painlev\'e表示との係数比較]
Chekhov--Mazzocco--Rubtsov \cite[Definition 1.1, (1.5)]{CMR} の量子化変数を
本稿と同じ $q$ とし、$\rho^2=q$、$\lambda=\rho(q^{-1}-q)$ と置く。
同文献の式に
\[
 (X_1,X_2,X_3)=(x,y,z),\qquad
 \epsilon_1=\epsilon_2=\epsilon_3=\lambda^{-1},\qquad
 (\Omega_1^0,\Omega_2^0,\Omega_3^0)=(a,b,c)/(q-1)
\]
を代入して各式に $\rho$ を掛けると、本稿の $t=1$ の三関係を得る。
例えば第一式は $xy-qyx-z-c=0$ になる。
三乗根では $\lambda(q-1)\ne0$ なので、この表示の特殊化は正則である。
ただし同文献のDefinition 1.1が仮定する非根条件は満たされない。
以下の三乗根でのPBW性と中心増大は、表示の一致に加えて直接証明する。
\end{remark}

\begin{proposition}\label{prop:pbw}
全係数で $B_{\mathbf a}$ はPBW基底 $t^\ell x^iy^jz^k$ を持ち、
各fiberはnoetherian quadratic Koszul CY AS4型14641の整域である。
低次数中心は
\[
Z(B_{\mathbf a})_1=kt,\qquad
Z(B_{\mathbf a})_3=k\Omega_{\mathbf a,h}\oplus kX_h\oplus kY_h\oplus kZ_h\oplus kt^3
\]
であり、非被約底を含む底変更と可換する。
\end{proposition}
\begin{proof}
$t=1$ の書換えは
$yx=q^2xy-q^2z-q^2c$、$zx=qxz+y+b$、$zy=q^2yz-q^2x-q^2a$。
唯一の重なり $zyx$ の二還元が全係数で一致する。
Rees斉次化とDiamond lemmaからPBW基底を得る。
活動次数の先頭環は三変数量子多項式環 $S$ の中央多項式拡大 $S[t]$ であり、
整域・noetherian性が従う。正則中央拡大のAS正則性 \cite{LBS} と二次Gr\"obner基底からAS4とKoszul性を得る。
Nakayama自己同型は $t$ を固定し $x_i\mapsto x_i+v_it$ と書ける。
第一関係の差は
$(1-q)t(v_yz+v_zy)+((1-q)v_yv_z-v_x)t^2=0$。
PBW独立性から全 $v_i=0$、よってCYである。

式\eqref{eq:Omega}と三つの $X_h$ の十二交換子は元の関係から零となる。
$S$ の中心単項式の条件は $i\equiv j\equiv k\pmod3$。
活動次数一・二には非定数中心がなく、三次には $x^3,y^3,z^3,xyz$ のみがある。
最高記号を引く帰納法で全低次数中心を得る。
非中心係数を消す際に使う $1-q,1-q^2$ は単元なので、底変更についても同じ議論が有効である。
\end{proof}

再開と検証に必要な全中心式を保持する。$t=1$ における元を $X,Y,Z,\Omega$ とし、
$f=(1-q)/3,e=-(2q+1)/9$ と置くと
\begin{align*}
\ell_a&=((2q+1)a^3+2(q-1)a+3qbc)/9,\quad \ell_b,\ell_c\text{ は巡回置換},\\
K&=((q+1)(a^2+b^2+c^2)+(q+2)abc)/3,\\
L&=((1-q)(a^2b^2+a^2c^2+b^2c^2)+(q+2)(a^2+b^2+c^2))/9.
\end{align*}
全中心は一本の関係
\begin{equation}\label{eq:center}
\begin{split}
0={}&\Omega^3+f\Omega^2-XYZ+e(X^2+Y^2+Z^2)\\
&-(aX+bY+cZ)\Omega+\ell_aX+\ell_bY+\ell_cZ+K\Omega+L
\end{split}
\end{equation}
で与えられる。元のPBW関係でこの恒等式を検証できる。
四生成元が先頭中心を生成することと、$\Omega$ に関するmonic三次式の剰余
$X^iY^jZ^k\Omega^r$ ($0\le r<3$) の相異なる最高記号により、これが全関係である。

$u=[\alpha:\beta:\gamma:\delta:\varepsilon]\in\PP^4$ に対して
\[
G_{\mathbf a,u}=\alpha\Omega_{\mathbf a,h}+\beta X_h+\gamma Y_h+\delta Z_h+\varepsilon t^3,
\qquad R_{\mathbf a,u}=B_{\mathbf a}/(G_{\mathbf a,u})
\]
とする。$\Delta=\alpha^3+27\beta\gamma\delta$、
$T_{\rm ell}=\{\beta\gamma\delta\Delta\ne0\}$ と置く。
境界は、平面三次
\[
E_u:\quad \beta x^3+\gamma y^3+\delta z^3+\alpha q^2xyz=0
\]
の $\phi=\operatorname{diag}(1,q,q^2)$ によるtwisted homogeneous coordinate ringである。
$T_{\rm ell}$ 上では三座標頂点を避ける滑らかな楕円曲線であり、$\phi$ は正確な位数三の平行移動となる。

\begin{assumption}[現行H3原稿から使う入力]\label{ass:recognition}
次の結果を既存の証明への明示的な依存として使う。
(1) $\mathbf a=0$ の切断のH3部分は非空開集合 $\mathcal U\subset\PP^4$ である。
(2) 被約有限型底上、PBW平坦な中央包絡・中央三次商族で、中央境界が滑らかな楕円THCR、
その平行移動が同一の有限位数を持つ場合、H3条件は開条件である。
(3) 中央滑らか楕円境界の場合、H3性は $R[t^{-1}]_0$ の有限大域次元と同値である。
(4) H3のpointed exact同値から反標準環と二次包絡が復元される。
証明の所在は研究ノートの固定参照表に記す。これらの認識結果を以下の代数計算のみから結論しない。
\end{assumption}

入力\ref{ass:recognition}よりH3開集合 $\mathcal V\subset\AA^3\times T_{\rm ell}$ を得る。
零切片は $\mathcal U_{\rm ell}=\mathcal U\cap T_{\rm ell}$ に一致するため、$\mathcal V$ は非空で七次元既約である。
全包絡同型は次数スカラーを除いて二符号変更と巡回置換だけであり、
\begin{equation}\label{eq:N}
\mathcal N=[\mathcal V/\mathfrak A_4]
\end{equation}
は七次元の滑らかな分離DM表示stackである。
実際、中心一次直線を保って $x_i\mapsto\sum_jM_{ij}x_j+v_it$ と書き、
境界のpoint graphを比較すると $M$ は対角・巡回型になる。
関係の $tx_j$ 係数が $v_i=0$ を強制し、残りから
$M=\lambda\operatorname{diag}(\epsilon_i)$、$\epsilon_i^2=1,\prod\epsilon_i=1$ を得る。
従って $a=b=c=0$ は余次元三の正則閉部分stackであり、固定包絡四次元族の非成分性が従う。

\section{固定位数三の全中央PBW正常形}\label{sec:normal}
\begin{theorem}\label{thm:normal}
$S=k\langle x,y,z\rangle/(yz-qzy,\mathrm{cyclic})$ の中央PBW拡大を
\[
r_i=t\sum_jA_{ij}x_j+b_it^2,\qquad [t,x_i]=0
\]
と書く。ここで $A$ は対称とする。
$A_{11}A_{22}A_{33}\ne0$ の開部分は、有限\'{e}taleな係数拡大の後、式\eqref{eq:B}へ正規化される。
従って七次元族は、この固定位数三・中央境界部分の指定開集合の全正常形である。
\end{theorem}
\begin{proof}
ポテンシャルを
\[
\Phi=xyz-qxzy-\tfrac12\sum_{i,j}A_{ij}x_ix_j-\sum_i b_ix_i
\]
と書く。二次以下の巡回語は対称式として扱える。
$x_i=X_i+v_i$ とすると、非対角二次係数を消す唯一の平行移動は
\[
v_1=A_{23}/(1-q),\quad v_2=A_{13}/(1-q),\quad v_3=A_{12}/(1-q)
\]
である。対角係数 $h_i=A_{ii}$ は変わらず、新しい定数は
\[
b'_i=b_i+h_iv_i+(1-q)v_jv_k\qquad(\{i,j,k\}=\{1,2,3\})
\]
となる。これはポテンシャルを直接展開した恒等式である。
次に $X_i=r_iY_i$、$t$ は固定とし、
\[
r_i^2=h_jh_k,\qquad r_1r_2r_3=h_1h_2h_3
\]
を選ぶ。単元の平方根を有限\'{e}taleに添加すれば可能で、選択肢は四つの偶符号である。
$h_ir_i=r_jr_k$ により対角一次係数はすべて1、新定数は $b'_i/(r_jr_k)$ となる。
残余群は前節の全同型計算と一致する。
\end{proof}

対称性は単なる便宜上の十分条件ではない。
巡回三次syzygyを一次補正へ作用させると
$\sum_{ij}A_{ij}[x_i,x_j]=0$ を得る。
$q\ne1$ の量子平面では三つの交換子が独立なので反対称部分は零である。
逆に対称 $A$ はKoszul PBW条件を満たす \cite{BG,EG}。
したがって本節は対称な表示だけを恣意的に選んだ分類ではない。
ただし中央境界を失う変形、非正規境界、対角係数零の閉部分はこの定理に含めない。

\section{三捩れから外へ出る全次数の変形}\label{sec:integrate}
\subsection{低次中心がないことによる持上げ}
\begin{lemma}[中央三次の正則持上げ]\label{lem:central-lift}
$A=k[[s]]$、$K=k((s))$ とする。
$V_A$ を濾過各次数で有限自由なPBW族、$\operatorname{gr}V_A=S_A$ とする。
$S_A$ に中央三次元 $g_A$ があり、$g_0$ の中心持上げ $G_0\in V_0$ が存在すると仮定する。
さらに $F_2Z(V_0)=k$、$\dim_K Z(S_K)_3=1$、
$V_K$ に次数三の非定数中心元が存在するとする。
このとき $g_A$ は $V_A$ の中心元へ持ち上がり、その定数項を固定すれば一意である。
指定した $G_0$ を特殊値に取れる。
\end{lemma}
\begin{proof}
$g_A$ の任意の持上げ $\widetilde g$ を選び、補正を
$F_2V_A/A$ の九つの基底で書く。
三生成元との交換子は有限自由加群間の線型方程式 $M(s)c=b(s)$ となる。
$F_2Z(V_0)=k$ なので $M(0)$ の階数は9であり、ある $9\times9$ minorは単元である。
この九行を使った解 $c$ は $A$ に属する。
一般fiberでは非定数中心元の最高記号が $g_K$ の非零倍であり、
規格化すれば全方程式の解が存在する。九行で解は一意なので、残る行も $K$ 上で零になる。
$A$ は整域であるから全行が $A$ 上で零である。
特殊fiberでも定数を除いた解は一意なので所望の $G_0$ へ戻る。
\end{proof}

\subsection{楕円曲線を固定し、平行移動を動かす}
有限な座標拡大の後、境界曲線を
\[
E_\lambda:\quad x^3+y^3+z^3-3\lambda xyz=0,\qquad \lambda^3\ne1
\]
とする。原点を $(1:-1:0)$、三捩れ点を $p_0=(1:-q:0)$ と取る。
\[
p_s=(1:b(s):s),\qquad 1+b(s)^3+s^3-3\lambda b(s)s=0,\qquad b(0)=-q
\]
を考える。$b$ に関する偏微分は $3q^2\ne0$ なので、$b(s)$ は一意の形式解を持ち、
実際に有限型曲線の\'{e}tale局所座標で定義される。
対応する三変数環は
\[
S_s=k\langle x,y,z\rangle/(yz+b(s)zy+sx^2,\ \mathrm{cyclic})
\]
である。

\begin{proposition}[全次数の包絡・切断変形]\label{prop:formal-integral}
七次元族の滑らかな楕円境界部分の任意の幾何学的点は、
境界の平行移動が三捩れから動く、全次数平坦なCY AS4包絡と正則中央三次商の形式族へ延長できる。
その一般fiberも滑らかな同じ楕円曲線を境界とする。
ここでは一般fiberのH3性はまだ結論しない。
\end{proposition}
\begin{proof}
上記のSklyanin表示は $p_s$ による平行移動の二次包絡である。
multilinear行列の行列式は
$b(s)s\{3\lambda xyz-(x^3+y^3+z^3)\}$。
一般fiberは非退化な楕円AS3、特殊fiberは既知の量子多項式環である。
全次数AS正則性は \cite{IS} と標準の楕円THCR認識を使う。
各次数のHilbert値は両fiberで等しいので、DVR上の有限加群の構造定理により $S_A$ は次数別に自由である。
THCR族への全射の三次核は階数一の束をなし、これを $kg_A$ と書く。
平行移動の場合のNakayama公式から $g_A$ は中心であり、$g_0$ は指定した境界三次式である。

ポテンシャルの二次・一次項を特殊fiberの値に固定してPBW変形を取る。
Etingof--Ginzburg \cite[Theorem 3.3.2(i)]{EG} と正則中央拡大の定理により、
Rees包絡は各fiberでCY AS4、全次数で平坦である。
同定理(ii)から一般fiberに非定数中央三次元が存在する。
一般平行移動は位数一・二・三ではなく、先頭環の低次数中心は定数と唯一の三次直線に限られる。
特殊fiberでは命題\ref{prop:pbw}により $F_2Z=k$。
補題\ref{lem:central-lift}を適用すると、指定した中央三次切断が正則に持ち上がる。
境界三次元の正則性と次数別乗法の単射により、商も全次数で平坦となる。
境界THCRが整域なので $t$-進最高記号から商の整域性も従う。
\end{proof}

\begin{lemma}[超越的中心値のH3性]\label{lem:generic-D}
前命題の包絡と中央三次持上げを中心値 $D$ から独立に定義し、
$D$ をそれらの定義体上超越的に取る。このとき一般fiberの切断 $G+D$ はH3である。
\end{lemma}
\begin{proof}
非斉次化 $V$ は有限大域次元を持つCY3環で、$G$ は非定数中心元である。
$G$ の最高記号が非零の正次数正則元なので $G$ は定義体 $F$ 上超越的である。
$F(D)$ 上の $V/(G+D)$ は、$D=-G$ と同定した中心局所化
$V\otimes_{F[G]}F(G)$ に等しい。
中心局所化は有限大域次元を保つ。
滑らかな中央楕円境界に対する入力\ref{ass:recognition}(3)を適用してH3を得る。
これは任意の数値 $D$ に関する開性や許容条件を主張する議論ではない。
\end{proof}

この構成は七次元族の幾何学的一般点から外へ出るH3一般fiberを与える。
境界データ $(E,\rho)$ の二次元、対称二次・一次ポテンシャルの九係数、
平行移動三次元、共通拡大一方向、中心値一方向を数えると
\[
2+9-3-1+1=8
\]
となる。
$1+b(s)$ が単元なので非対角二次項は平行移動で一意に消せる。
一般の境界の射影自己同型群は有限であり、対角二次係数の一つを共通拡大で規格化すれば、
この計算は単なる冗長パラメータ数でなく、八次元の一般表示族の次元を与える。
全intrinsic族への射を大域的に定義したこととは区別する。

\section{滑らかな中心でのHochschild変形空間}\label{sec:hh}
\begin{assumption}[本節の幾何学的適用範囲]\label{ass:order}
$C$ は滑らかな射影三次曲面、$D\in|-K_C|$ は滑らかな楕円曲線とする。
$\Lambda$ は $C$ 上の有限orderで、$C\setminus D$ ではAzumaya、
$D$ 上では\'{e}tale局所的な完備化が
\[
k[[s,w]]\rtimes\mu_3,\qquad \zeta(s,w)=(s,\zeta w)
\]
とMorita同値であるとする。
七次元族の滑らかな中心を持つ開部分は、既存の中心・境界order定理によりこの範囲に入る。
原点切片には第\ref{sec:exact}節の厳密な滑らかさの証明書を持つ点がある。
\end{assumption}

\begin{lemma}[局所Hochschild層]\label{lem:local-HH}
仮定\ref{ass:order}の下で、$C$ 上の局所Hochschild cohomology層は
\[
\mathcal{HH}^0=\OO_C,\quad
\mathcal{HH}^1=T_C(-\log D),\quad
\mathcal{HH}^2=\bigwedge^2T_C(-\log D),\quad
\mathcal{HH}^{>2}=0
\]
である。
\end{lemma}
\begin{proof}
非分岐部分は局所Morita不変性と通常のHKRである。
分岐点では $A=k[[s,w]]$ の二変数Koszul分解を使い、
$\Ext_{A^e}(A,A\rtimes\mu_3)^{\mu_3}$ を計算する。
恒等元の部分は $\bigwedge T_A$ の不変部分であり、その基底は
\[
1;\qquad \partial_s,\ w\partial_w;\qquad w\partial_s\wedge\partial_w.
\]
$z=w^3$ とすると $w\partial_w=3z\partial_z$ なので、これらは
$\OO_C,T_C(-\log D),\wedge^2T_C(-\log D)$ に対応する。

非恒等元の固定部分は $w=0$ の直線である。
Koszul分解の $w$ 方向を計算すると、寄与は法線表現の行列式を伴う。
この一次元表現への $\mu_3$ 作用は非自明なので、次数一・二の両方で不変部分は零である。
これは \cite[Theorem 1.1]{NS} の固定部分公式の直接計算でもある。
この局所計算はBrauer捩れを消す大域的同値を必要としない。
局所Morita同値の変更は、恒等部分の中心に由来するpolyvector層に恒等に作用する。
それらの同定は非分岐稠密開上で通常の中心の同定に一致し、局所自由層への延長が一意であるから貼り合う。
非恒等部分は局所的に零なので、Brauer捩れから新しい層が生じることもない。
\end{proof}

\begin{theorem}[八次元の無障害な圏変形]\label{thm:HH8}
$\mathcal C=\Perf(C,\Lambda)$ とする。仮定\ref{ass:order}の下で
\[
\HH^0(\mathcal C)=k,\qquad \HH^1(\mathcal C)=0,\qquad
\dim\HH^2(\mathcal C)=8,\qquad \HH^n(\mathcal C)=0\ (n\ge3).
\]
さらに標準的な完全列
\[
0\longrightarrow H^1(C,T_C(-\log D))\longrightarrow\HH^2(\mathcal C)
\longrightarrow H^0(C,\OO_C)\longrightarrow0
\]
があり、左右の次元は7と1である。
従って通常のHochschild複体で制御される圏の形式変形は無障害で、
次数スカラーの高次自己同型を別にすれば、そのminiversal底は八次元形式円板である。
\end{theorem}
\begin{proof}
滑らかな三次曲面について $H^{>0}(\OO_C)=0$、
$h^0(T_C)=0,h^1(T_C)=4,h^2(T_C)=0$。
後二つは $C\subset\PP^3$ の法束列とEuler列、
$h^0(\OO_C(3))=19$、$h^0(T_{\PP^3}|_C)=15$ から得る。
$h^0(T_C)=0$ は、連結自己同型が27本の直線を固定し、
六本を縮約した $\PP^2$ 上で一般位置の六点を固定することからも従う。

対数接束の完全列
\[
0\longrightarrow T_C(-\log D)\longrightarrow T_C
\longrightarrow\OO_D(D)\longrightarrow0
\]
を使う。楕円曲線上の $\OO_D(D)$ の次数は3なので $h^0=3,h^1=0$。
従って対数接束のcohomologyの次元は順に $0,7,0$ である。
また $D=-K_C$ から $\wedge^2T_C(-\log D)=\OO_C$。

局所から大域へのHochschild spectral sequenceを補題\ref{lem:local-HH}に適用する。
非零項は $H^0(\OO_C)$、$H^1(T_C(-\log D))$、
$H^0(\wedge^2T_C(-\log D))$ に限られ、すべての微分は次数の理由で零である。
一般的な大域HKR形式性を別途仮定せず、主張と完全列を得る。
通常の変形理論 \cite{LVdB} では接空間が $\HH^2$、小拡大に対する障害が $\HH^3$ に属する。
後者の消滅から形式的滑らかさ、前者の有限次元性から八次元のminiversal底を得る。
\end{proof}

\begin{corollary}
構造対象 $\OO$ がexceptionalである場合、その持上げの障害と選択の曖昧さ \cite{LowenObject} は
それぞれ $\Ext^2(\OO,\OO)$ と $\Ext^1(\OO,\OO)$ に属し、いずれも零である。
従って構造対象を持つ圏の形式変形でも、接空間と障害の上記の値は変わらない。
\end{corollary}

\begin{remark}
この定理は、全H3変形が自動的に同じ14641二次包絡を持つこと、
逆Serre作用が任意の変形heartを保つこと、全次数の相対反標準環が再構成できることを証明していない。
圏の八次元形式変形と、全H3スタックの八次元開部分を同一視するには、これらの比較が必要である。
\end{remark}

\section{厳密計算と具体的な横断方向}\label{sec:exact}
非斉次化 $t=1$、$a=b=c=0$ とし、
\[
\Phi_0=xyz-qxzy-\tfrac12(x^2+y^2+z^2),\qquad
G_0=2\Omega+X+2Y+3Z+5
\]
を取る。$\Delta=170$、$\beta\gamma\delta=6$。
式\eqref{eq:center}に $\Omega=-(X+2Y+3Z+5)/2$ を代入した三次式を $F$ とすると、
$\mathbb Q(q)$ 上で $(F,F_X,F_Y,F_Z)$ のGr\"obner基底は $(1)$ である。
既存の境界正則性定理と合わせ、この点は滑らかな中心と楕円境界を持つH3点である。

次数一から三の巡回語は20個、次数三以下のPBW語も20個ある。
それぞれ $\dot\Phi,\dot G$ の係数を未知数とし、
変形した関係式の下で $[G_0+\epsilon\dot G,x_i]=0\pmod{\epsilon^2}$ を一次化する。
全105行40列の行列 $M$ を $\mathbb Q(q)$ 上で計算すると
\[
\operatorname{rank}M=18,\quad\dim\ker M=22.
\]
12個のaffine座標変更、ポテンシャルのスカラー、切断のスカラーの14列を明示すると、
すべて核に入り階数14である。
既知の七次元族を加えたspanの階数は21であり、次の方向を加えると22になる：
\[
\dot\Phi=x^3+2y^3+3z^3+2xyz,
\]
\begin{align*}
\dot G={}&(22q+6)x+(5q+9)y+(18q+18)z\\
&-\tfrac43(q+1)(x^2+z^2)-\tfrac43(q+4)y^2+6xy+9q\,xz+18yz.
\end{align*}
従って\emph{このポテンシャル・中央三次incidence}の接空間商は正確に八次元である。
これは全H3変形の被覆を行列の階数だけから推論したものではない。

全次数の明示例として、別の許容候補切断
$2\Omega+X+8Y+27Z+D$ に対し
\begin{equation}\label{eq:explicit-integral}
\Phi_u=(1+2u)xyz-qxzy+u(x^3+8y^3+27z^3)-\tfrac12(x^2+y^2+z^2)
\end{equation}
を取る。Etingof--GinzburgのAppendix 9.1の中心式を
\[
X_{\rm R}=x,\quad Y_{\rm R}=2y,\quad Z_{\rm R}=3z,
\quad q_{\rm R}=q/(1+2u),\quad t_{\rm R}=-18u/(1+2u)
\]
で変換すると、$u=0$ で正則な有理中心三次元を得る。
$x^3$ 係数を1、定数を $D$ と規格化した特殊値は正確に
$2\Omega+X+8Y+27Z+D$ である。
全有理係数は小さいJSONに保持し、原関係からの三つの交換子を
$\mathbb Q(q)(u)$ 上で独立に零へ還元した。
この検算は原典の中心性を前提としない。
$D=5$ と固定した一般 $u$ のH3性をこの計算だけで主張しない。
一般中心値 $D$ については補題\ref{lem:generic-D}が適用される。

\section{Artin底上の反標準切断と変形の単射性}\label{sec:artin-reconstruction}
この節では、一次元のDVR上の積分とは別に、任意の局所Artin底上での再構成を扱う。
全圏変形がH3になるとは結論しない。解決するのは、切断環の平坦性・有限表示と、
適合したenhancementを持つAbel変形から圏変形への\emph{同型類上の単射性}である。

\subsection{反標準列に対する消滅と既知の再構成定理}
$R_0=B_0/(G_0)$ を本稿のH3点の飽和反標準環とし、
$\mathcal A_0=\qgr R_0$、$E_i=\pi R_0(i)$、
$\mathcal G_0=\operatorname{Ind}(\mathcal A_0)=\operatorname{QGr}R_0$ と置く。
$R_0$ はAS Cohen--Macaulay、AS Gorenstein次元3・parameter1である。
標準的な局所cohomologyとtails cohomologyの比較および双対性から
\begin{equation}\label{eq:artin-forward-vanishing}
 \Hom(E_i,E_j)=(R_0)_{j-i},\qquad
 \Ext^p(E_i,E_j)=0\quad (p>0,\ j\ge i)
\end{equation}
となる。実際、第一cohomologyはCohen--Macaulay性で消え、
第二cohomologyは $(R_0)_{i-j-1}^{*}=0$ となる。
特に全 $E_i$ はexceptionalであり、$\{E_i\}_{i\in\mathbb Z}$ はAZ ampleな列である。

正方向だけを残した有向 $\mathbb Z$-代数を
\[
 \mathfrak a_0(i,j)=
 \begin{cases}\Hom(E_i,E_j)&i\le j,\\0&i>j\end{cases}
\]
と書く。右加群の規約では添字を $i\mapsto-i$ と変えてもよい。
\cite[Theorem 4.5]{DeDekenLowen}を適用すると、局所Artin $k$-代数 $A$ に対し
\begin{equation}\label{eq:DDL-equivalence}
 \operatorname{Def}_{\rm lin}(\mathfrak a_0)(A)
 \simeq \operatorname{Def}_{\rm ab}(\mathcal G_0)(A)
\end{equation}
を得る。ここでAbel変形のflat性はLowen--Van den Berghの意味である。
必要なExt消滅は\eqref{eq:artin-forward-vanishing}であり、
有限次元 $k$-空間 $X$ をtensorした場合も有限直和だから同じである。
$\mathbb Z$-生成性はAZ再構成から従う。
変形されたAbel圏における $E_i$ のflatな持上げを $E_{i,A}$ とすると、
変形 $\mathfrak a_A$ はそれらの正方向Homから構成される。
noetherian性などの有限性の持上げは同論文Proposition 4.6による。
これは既知の定理の適用であって、新しい一般的なAbel変形理論ではない。
また、ここでは $\mathfrak a_A$ の一周期性を仮定も結論もしていない。

\subsection{enhancementを持つAbel変形は圏変形に埋め込まれる}
\begin{definition}\label{def:compatible-abelian}
本節でいう\emph{適合したenhanced Abel変形}は、\eqref{eq:DDL-equivalence}のflat Abel変形と、
そのperfect部分の $A$-線型dg enhancementで、次を満たすものである。
$E_{i,A}$ はperfect対象として実現され、特殊fiberは指定された $E_i$ である。
perfect対象間のderived Homは $A$ 上perfectであり、その形成が導来底変更と両立する。
射には特殊fiberで恒等となるmarkingを課す。
この適合性を、任意の抽象dg変形にheartが自動的に付くという主張に置き換えない。
\end{definition}

\begin{theorem}[全Artin底での同型類の単射性]\label{thm:artin-injective}
二つの適合したenhanced Abel変形のdg圏が、特殊fiberのmarkingを保つ
$A$-線型Morita同値で結ばれているとする。
このとき二つのAbel変形は、構造対象を保つmarked exact同値によって同型である。
従って忘却写像は同型類の集合上単射であり、その接写像も単射である。
\end{theorem}
\begin{proof}
Morita同値を $F$ とする。任意の整数 $i$ に対して
$F(E_{i,A})$ と $E'_{i,A}$ は、同じ $E_i$ の、同じdg変形内のperfectな持上げである。
小拡大 $A\twoheadrightarrow\bar A$ の核 $J$ に対する対象持上げの差は
$\Ext^1(E_i,E_i)\otimes_kJ$ に、存在障害は $\Ext^2(E_i,E_i)\otimes_kJ$ に属する
\cite{LowenObject}。いずれも零である。
Artin長による帰納法で、markingを持ち上げる同型
$F(E_{i,A})\simeq E'_{i,A}$ を各 $i$ について選べる。
選択が一意な射であるとは主張しない。対象のスカラー自己同型は残る。

これらの同型と $F$ の合成保存性により、正方向の全Homと積を保つ同型
$\mathfrak a_A\simeq\mathfrak a'_A$ を得る。
両辺のHomはheartのHomでもあり、perfect enhancement内の次数零Homでもある。
従って\eqref{eq:DDL-equivalence}を逆に適用すると、Abel変形のexact同値を得る。
添字0の対象も保たれるのでpointedである。
$A=k[\epsilon]/(\epsilon^2)$ を取れば接写像の単射性が従う。
\end{proof}

\begin{remark}
この定理は与えられた $F$ 自体のt-exact性や、dg自己同値群oidへの充満忠実性を主張しない。
必要なのは二つのAbel変形の同型類を区別できることである。
また、無限列を有限な完全例外列に取り替えていない。
三次曲面の反標準列は通常、有限長の例外列とは異なるので、
\cite{DeDekenLowen}のgeometric helixに関する第5節はこの証明では使用しない。
\end{remark}

\begin{corollary}\label{cor:artin-tangent-bound}
定理\ref{thm:HH8}の一般点において、適合したenhanced Abel変形として定義される
intrinsic H3変形の接空間は8次元以下である。
従って定理\ref{thm:conditional-component}の条件(2)は、この相対変形問題では
追加仮定ではなく定理\ref{thm:artin-injective}の帰結となる。
条件(1)の局所有限型代数性、条件(3)の相対H3族の実現、全境界被覆は依然として必要である。
\end{corollary}
\begin{proof}
H3の構造対象はexceptionalで、その持上げは同型類として一意である。
前定理を圏の変形の接空間 $\HH^2$ と比較し、定理\ref{thm:HH8}を用いる。
\end{proof}

\subsection{相対反標準切断環の有限自由性と六二次・一三次表示}
\begin{lemma}[perfectな複体の特殊fiberからの認識]\label{lem:artin-perfect}
$A$ を局所Artin環、剰余体を $k$ とし、$P\in\Perf(A)$ とする。
$P\otimes_A^{\mathbf L}k\simeq k^r[0]$ なら $P\simeq A^r[0]$ である。
この同定によりcohomologyの形成は任意の $A$-代数への底変更と両立する。
\end{lemma}
\begin{proof}
有界有限自由複体から、微分の単元成分を持つ可縮な二項直和因子を順に除く。
残る最小複体の微分は極大イデアルに入るので、$k$ とのtensor後は零になる。
特殊fiberのcohomologyが次数零の $k^r$ のみであることから、
この最小複体は次数零の $A^r$ だけとなる。
これは\cite[Tag 0BDK]{StacksPerfect}のArtin局所版でもある。
\end{proof}

\begin{theorem}[全Artin底での反標準環の平坦な有限表示]\label{thm:artin-section-ring}
$\mathcal D_A$ をsmooth properな $A$-線型dg圏とし、
特殊fiberを $\Perf(\mathcal A_0)$ と同定する。
相対Serre自己同値 $S_A$ とその逆が特殊fiberへの底変更に両立し、
$E_A$ が $\OO$ のperfectな持上げであるとする。
$\tau_A=S_A^{-1}[2]$ と置く。このとき
\[
 R_A=\bigoplus_{n\ge0}H^0\operatorname{RHom}_{\mathcal D_A}(E_A,\tau_A^n E_A)
\]
は、自己同値による通常のorbit積について、次数別に有限自由なconnected graded $A$-代数である。
その各次数のrankは $(3n^2+3n+2)/2$、特殊fiberは $R_0$、次数別の形成は任意のArtin底変更と両立する。
さらに $R_A$ は左右noetherianで、一次生成であり、
\begin{equation}\label{eq:artin-six-one}
 R_A\simeq T_A(A^4)/(r_1,\ldots,r_6,g),\qquad \deg r_i=2,\quad\deg g=3
\end{equation}
という表示を持つ。六つの二次関係は関係空間の $A$-基底に取れ、$g$ は特殊fiberの唯一の最小三次関係を持ち上げる。
\end{theorem}
\begin{proof}
各derived Homはproper性によりperfectであり、特殊fiberは
\eqref{eq:artin-forward-vanishing}によって $(R_0)_n[0]$ である。
補題\ref{lem:artin-perfect}を適用して有限自由性と底変更を得る。
$n=0$ では単位写像 $A\to R_{A,0}$ が剰余体上同型なので同型である。
$\tau_A$ の反復を用いた積は結合的である。
各次数の乗法 $R_{A,1}^{\otimes n}\to R_{A,n}$ の余核は有限生成で、
特殊fiberで零なのでNakayama補題から零となる。従って一次生成である。

$T=T_A(R_{A,1})$、$I=\ker(T\to R_A)$ とする。
$0\to I_n\to T_n\to R_{A,n}\to0$ は $A$-加群として分裂するため、
$I_n$ は有限自由で、その特殊fiberは $(I_0)_n$ と一致する。
従って $I_2$ のrankは $16-10=6$ である。
$I_2$ の基底 $r_i$ と、特殊fiberの最小三次関係を持ち上げる $g\in I_3$ を選ぶ。
これらで生成される両側イデアルを $J$ とする。
$R_0=B_0/(G_0)$ の全関係は二次関係と一つの三次関係で生成されるので、
各次数で $J_n$ の $I_n\otimes_Ak$ への像は全体である。
従って有限生成加群 $I_n/J_n$ の特殊fiberは零であり、再びNakayama補題から $I_n=J_n$。
これが\eqref{eq:artin-six-one}である。

$\mathfrak m^N=0$ となる $N$ を取り、
$R_A$ の有限な $\mathfrak m$-進filtrationを取る。
$A$-flat性から各商は
$R_0\otimes_k(\mathfrak m^j/\mathfrak m^{j+1})$ であり、
左右の $R_A$-加群として有限個のnoetherianな $R_0$ の直和になる。
有限拡張により $R_A$ 自身が左右noetherianである。
\end{proof}

\begin{remark}[残っている再構成条件]\label{rem:artin-gap}
定理\ref{thm:artin-section-ring}ではheartを仮定しないので、$\tau_A$ のt-exact性は結論しない。
また $B_A=T_A(A^4)/(r_1,\ldots,r_6)$ の平坦性・AS4性は、$R_A$ の平坦性だけからは導いていない。
$B_A\to R_A$ の核が両側イデアル $(g)$ であることと、
$g$ が正則正規でその核が左右の主イデアルであることは異なる。
指定potential chart内では次節の定理\ref{thm:artin-potential-envelope}がこの包絡条件を解決する。
一般の全圏変形がそのchartに入ることはなお別の認識問題である。
本節の主張は通常の証明であり、以下の有限計算証明書により代替されるものではない。
\end{remark}

\subsection{二次包絡の三次平坦性を支配する四係数}
\label{sec:quadratic-four-obstructions}
前節の再構成に残る条件のうち、三次の条件を有限行列として分離できる。
以下では全次数平坦性を三次平坦性から無条件に推論しない。

\begin{proposition}\label{prop:four-cubic-obstructions}
$A$ を剰余体 $k$ を持つ可換局所Artin環とし、$R_A$ を
定理\ref{thm:artin-section-ring}の平坦な切断環とする。
$V=R_{A,1}$、$I_n=\ker(V^{\otimes n}\to R_{A,n})$、
$B_A=T_A(V)/(I_2)$ と置く。
特殊fiberは $R_0=B_0/(G_0)$ で、$B_0$ は14641型Koszul AS4、
$G_0$ は正則三次元とする。
このとき、基底の選択後、四つの係数 $o_1,\ldots,o_4\in\mathfrak m_A$ があり、
\[
 (B_A)_3\simeq A^{19}\oplus A/(o_1,\ldots,o_4).
\]
従って $(B_A)_3$ が $A$-flatであることと、四係数が全て零であることは同値である。
イデアル $(o_1,\ldots,o_4)$ は基底によらず
\[
 \operatorname{Fitt}_0\bigl(I_3/(V I_2+I_2V)\bigr)
\]
である。特にこの障害イデアルの形成は、この有限表示の任意の底変更と両立する。
\end{proposition}
\begin{proof}
$V,I_2,I_3$ はそれぞれrank $4,6,45$ の有限自由加群である。
二次関係挿入写像
\[
 d:(V\otimes I_2)\oplus(I_2\otimes V)\longrightarrow I_3
\]
は $45\times48$ 行列で表示される。
特殊fiberで像の次元は $64-\dim(B_0)_3=44$ である。
従ってある $44\times44$ 小行列式は単元であり、可逆な行・列操作で
\[
 d\ \sim\
 \begin{pmatrix}I_{44}&0\\0&(o_1\ o_2\ o_3\ o_4)\end{pmatrix}
\]
とできる。剰余体上のrankが44なので $o_i\in\mathfrak m_A$ である。
これは $I_3/\operatorname{im}d\simeq A/(o_1,\ldots,o_4)$ とFittingイデアルの式を与える。
$R_{A,3}\simeq A^{19}$ なので
\[
 0\longrightarrow I_3/\operatorname{im}d\longrightarrow(B_A)_3
 \longrightarrow R_{A,3}\longrightarrow0
\]
は $A$-加群として分裂する。
最後に、局所Artin環上の有限flat加群は自由であり、
$A/J$ がflatならその剰余体上のrankは1なので $A/J\simeq A$、従って $J=0$ である。
逆向きは明らかである。
\end{proof}

\begin{lemma}[全次数平坦性の後の三次核]\label{lem:flat-envelope-principal}
上の設定で、さらに $B_A$ が次数別に $A$-flatであると仮定する。
$g\in\ker(B_A\to R_A)_3$ を $G_0$ の持上げとすると、
\[
 \ker(B_A\to R_A)=B_Ag=gB_A
\]
であり、$g$ は左右正則正規である。
従って全次数平坦性が得られた後には、三次核の正則正規主イデアル性を
別の無限個の方程式として課す必要はない。
\end{lemma}
\begin{proof}
$K_n=\ker((B_A)_n\to R_{A,n})$ と置く。
両辺が有限自由で右辺への写像は全射なので $K_n$ も有限自由で、底変更に適合する。
そのrankは $\dim(B_0)_n-\dim(R_0)_n=\dim(B_0)_{n-3}$ である。
左右の乗法写像
\[
 (B_A)_{n-3}\xrightarrow{\ b\mapsto bg\ } K_n,
 \qquad
 (B_A)_{n-3}\xrightarrow{\ b\mapsto gb\ } K_n
\]
は特殊fiberで $G_0$ の正則性と正規性によって同型である。
同rank有限自由加群間の写像なので、行列式が単元となり双方とも同型である。
全次数でこれを適用すれば主張を得る。
\end{proof}

\begin{proposition}[全20方向の同時一次包絡は全次数で平坦]
\label{prop:square-zero-potential-pbw}
$k$ は標数零で原始三乗根 $q$ を含むとし、
$A=k[e_1,\ldots,e_{20}]/(e_1,\ldots,e_{20})^2$ とする。
量子so3点のpotential $\Phi_0$ に、次数1--3の20個の巡回語をそれぞれ係数 $e_j$ で加える。
中央変数 $t$ で斉次化した六つの二次関係の包絡 $B_A$ は、
\[
 \{x^ay^bz^ct^d:a,b,c,d\ge0\}
\]
を $A$-基底に持つ。従って全次数で有限自由であり、
相対Koszul解消のrank列は $(1,4,6,4,1)$、双対cohomologyは最高次数の $A(4)$ のみである。
これはArtin底上の相対AS4性であり、非被約底環を含む $B_A$ 自身の有限大域次元を意味しない。

さらにこの包絡中に中央三次元 $g$ があり、特殊fiberが正則元 $G_0$ なら、
$g$ は左右正則で $B_A/(g)$ は次数別に有限自由である。
特に既存の一次potential・中央三次incidenceに属する全方向は、
包絡と切断環の平坦性に関しては実際の同時一次変形として実現される。
全intrinsic H3の変形の被覆は別問題である。
\end{proposition}
\begin{proof}
単なる四つのrank条件の外挿ではなく、平方零係数による書換えの停止性を使う。
六つの逆順対の書換えを、定数部分と平方零部分に分ける。
$e_j$ を既に含む項には定数部分だけを適用する。二回目の変形項は零になる。
斉次次数を固定し、まず平方零係数を含む語を係数を含まない語より小さく順序付ける。
同じ係数層では、$x,y,z$ の個数、次に $x<y<z<t$ の辞書式順序を使う。
定数部分の書換えはこの順序を狭義に下げ、変形項は係数層を下げる。
固定次数の語は有限個なので、全ての書換えは停止する。

左辺は六つの相異なる長さ2の逆順対であり、包含曖昧性はない。
重なる曖昧性は $zyx,tyx,tzx,tzy$ の四つだけである。
その定数部分は既知の量子so3 PBW恒等式、一次部分は
\texttt{verify\_cubic\_obstruction.py}の全80恒等式で零になる。
二次以上の係数積は底環で零である。
離れた二箇所の書換えは可換なので、全局所曖昧性が解消する。
停止する線型書換えのDiamondの議論により、既約語は主張した $A$-基底となる。

相対Koszul性も全次数flat性から確認できる。
augmentation $B_A\to A$ の被約bar複体の内部次数 $n$ 部分は有界有限自由 $A$-複体である。
特殊fiberのcohomologyは $B_0$ のKoszul性によりhomological次数 $n$ のみにあり、
そのrankは $n=0,1,2,3,4$ で $1,4,6,4,1$、それ以外で零である。
補題\ref{lem:artin-perfect}の議論を各内部次数に適用すると、相対Torも同じ対角に有限自由として集中する。
この基底を順に持ち上げて極小相対自由解消を構成すれば、
$B_A,B_A(-1)^4,B_A(-2)^6,B_A(-3)^4,B_A(-4)$ の長さ4の解消を得る。
その双対複体の各内部次数も有界有限自由である。
特殊fiberのAS Gorenstein性と同じ最小複体の議論から、双対cohomologyは最高次数の $A(4)$ のみとなる。

最後に各次数の $g$ による乗法は、有限自由加群間の写像で特殊fiberが単射である。
従って適切な最大小行列式が単元で、写像は分裂単射、余核は有限自由である。
左右両方に適用して正則性と商のflat性を得る。
\end{proof}

\begin{proposition}[巡回potential上の四syzygyは全Artin底で持ち上がる]
\label{prop:universal-four-syzygies}
量子so3点を特殊fiberとする任意の局所Artin底上の巡回potential変形について、
斉次二次包絡の三次部分は自由rank20である。平方零の仮定は不要である。
\end{proposition}
\begin{proof}
$\widetilde\Phi$ を四変数で三次に斉次化し、
$r_i=\partial_{x_i}\widetilde\Phi$、$r_t=\partial_t\widetilde\Phi$、
$c_i=[t,x_i]$ とする。六つの関係は $c_i,r_i$ で取れる。
$t$ を末尾へ移した表示との差は $c_i$ の線型結合なので、同じ二次包絡である。
二次式 $f=\sum a_{uv}uv$ に対し
\[
 D_t(f)=\sum a_{uv}(c_u v+u c_v),\qquad c_t=0
\]
と書く。自由代数で $D_t(f)=[t,f]$ なので、三つの恒等式
\[
 t r_i-r_i t-D_t(r_i)=0
\]
を得る。さらに巡回微分のEuler恒等式から
\[
 \sum_i(x_i r_i-r_i x_i)+D_t(r_t)=0
\]
を得る。これらは六関係の四つの線型syzygyであり、係数に打切りを課していない。

特殊fiberの64行48列の挿入行列はrank44で、上の四syzygyはその核の基底である。
後者は、$r_i$ の左係数に $t$ が現れる三本と、$x_i$ が現れる一本の独立性からも分かる。
\texttt{verify\_universal\_cubic\_syzygies.py}は全20係数の多項式環上で四恒等式を確認し、
特殊fiberの二行列のrank44と4も厳密確認する。
44単元minorによって挿入行列を $I_{44}$ と残り20行4列のSchur補行列へ変形する。
四syzygyの残り4成分は特殊fiberで基底なので可逆な4次行列をなし、
正確なsyzygy恒等式からこのSchur補行列は零である。
従って挿入行列の余核は $A^{20}$ である。
\end{proof}

\begin{theorem}[巡回potential chartの全Artin底での包絡平坦性]
\label{thm:artin-potential-envelope}
$A$ を剰余体 $k$ を持つ局所Artin $k$-代数とする。
$k$ は標数零で原始三乗根 $q$ を含むとし、
\[
 \Phi_A\in A\langle x,y,z\rangle_{\rm cyc}^{\le3},\qquad
 \Phi_A\bmod\mathfrak m_A=
 xyz-qxzy-\tfrac12(x^2+y^2+z^2)
\]
とする。三次に斉次化した $\widehat\Phi_A\in A[t]\langle x,y,z\rangle_{\rm cyc}$ に対し
\[
 B_A=A[t]\langle x,y,z\rangle/
 (\partial_x\widehat\Phi_A,\partial_y\widehat\Phi_A,
 \partial_z\widehat\Phi_A),\qquad \deg t=\deg x=\deg y=\deg z=1
\]
と置く。この二次包絡は全次数で有限自由、rankは $\binom{n+3}{3}$ であり、
形成は任意のArtin底変更と両立する。
$t$ は中央正則で、$B_A$ は左右noetherian、相対Koszul型14641・相対AS Gorensteinである。
中央三次元が与えられ特殊fiberで正則なら、その商も次数別有限自由となる。
全intrinsic変形がこのchartを通ること、全中央三次元の持上げ、H3族の大域的実現は本定理に含まない。
\end{theorem}
\begin{proof}
単に三次のflat性から高次を推論せず、potentialのdg複体を使う。
$A[t]$ 上の自由dg代数 $\mathcal D_A$ を、生成元
\[
 x_i:(\text{内部次数},\text{homological次数})=(1,0),\quad
 \xi_i:(2,1),\quad \zeta:(3,2)
\]
と微分
\[
 d x_i=d t=0,\qquad d\xi_i=\partial_{x_i}\widehat\Phi_A,
 \qquad d\zeta=\sum_i[x_i,\xi_i]
\]
で定める。$t$ は係数環の中央変数であるから、巡回微分恒等式により $d^2=0$ であり、
$H_0(\mathcal D_A)=B_A$ となる。
各内部次数 $n$ の部分は有界有限自由 $A$-複体である。

まず特殊fiber $\mathcal D_0$ を考える。
$t=0$ での商は斉次potential $xyz-qxzy$ の標準potential dg代数である。
このpotentialはCY3量子多項式環を定めるため、正のhomological次数でacyclicである。
ここで使う外部入力は \cite[\S7.1, (7.1.4)]{EG}
（そこで明示されたGinzburg Theorem5.3.1）である。
同箇所の定義・仮定を一次資料で確認しており、一般の非斉次potentialに逆向きを適用してはいない。
$\mathbb Q(q)$ 上の複体の正確性は複素数体への忠実平坦拡大から降下し、任意の $k$ へ拡大できる。

複体の短完全列
\[
 0\longrightarrow\mathcal D_0(-1)\xrightarrow{t}\mathcal D_0
 \longrightarrow\mathcal D_0/(t)\longrightarrow0
\]
の長完全列から、$j>0$ について $H_j(\mathcal D_0)=tH_j(\mathcal D_0)$ となる。
内部次数は下に有界なので、この等式は $H_j(\mathcal D_0)=0$ を意味する。
同時に $t$ は $H_0$ 上単射であり、$H_0/(t)$ は量子三変数多項式環となる。
従って $\dim H_0(\mathcal D_0)_n=\binom{n+3}{3}$ である。

任意の $n$ に対し、有限自由複体 $(\mathcal D_A)_n$ の特殊fiberのhomologyは次数零のみにある。
補題\ref{lem:artin-perfect}から、この複体は $A^{\binom{n+3}{3}}[0]$ と準同型である。
これが全次数flat性と底変更を同時に与える。底変更先が非flatでも、perfect複体から得た同定なので問題はない。
$t$ の正則性は次数ごとの有限自由加群間の写像を特殊fiberで判定して従う。
左右noetherian性は有限 $\mathfrak m_A$-進filtrationの商が $B_0$ の有限直和であることから従う。
相対Koszul解消と双対性は命題\ref{prop:square-zero-potential-pbw}のbar複体による議論を適用する。
その議論で平方零性を使ったのはflat性を得る最初の書換えだけであり、flat性を得た後は任意のArtin底で有効である。
最後の中央三次商も、同命題の分裂単射の議論で従う。
\end{proof}

\begin{proposition}[八方向の同時二次中心切断]
\label{prop:eight-second-order}
第\ref{sec:exact}節の $a=b=c=0$、切断係数 $(2,1,2,3,5)$ の点で、
\[
 A_2=k[h_0,\ldots,h_7]/(h_0,\ldots,h_7)^3
\]
上の巡回potentialと中央三次切断の同時変形が存在する。
その一次の八方向は座標変更と二つのスカラーを除いた一次incidenceの全体を張る。
包絡と三次商は全次数で $A_2$-flatである。
この二次jetの存在は、任意次数の同時形式積分や全intrinsic変形の被覆を意味しない。
\end{proposition}
\begin{proof}
\texttt{verify\_second\_order\_incidence.py}が有理係数を $\mathbb Q(q)$ 上で再生する。
既知七方向を14個の座標方向に順に加え、rankを増す七本を選び、
既存の横断方向を加える。全rankは22である。
この八方向について中央交換子の二次項を105行36列行列 $Q$ にまとめる。
元の一次行列を $M$ とすると
\[
 \operatorname{rank}M=\operatorname{rank}(M\mid-Q)=18
\]
である。従って $MH=-Q$ を厳密に解け、$H$ の40行36列が
potentialの20係数と三次元の20係数への二次補正を与える。

二次乗法では、逆順三対の係数行列も変化することに注意する。
それを $A_0+A_1+A_2$、既約語の係数を $C_0+C_1+C_2$ と書き、書換えを
\[
 B_0=-A_0^{-1}C_0,\quad
 B_1=-A_0^{-1}(C_1+A_1B_0),\quad
 B_2=-A_0^{-1}(C_2+A_1B_1+A_2B_0)
\]
で二次まで正しく展開する。
一次で正規化した関係をそのまま二次へ延長していない。
補正後の元関係からこの書換えを再構成し、三つの中央交換子の次数0,1,2を
全て再計算して零であることを確認する。
平方零係数のときと同じく、係数次数が上がるたび残りの許容次数が下がるため、この有限展開は停止する。
定理\ref{thm:artin-potential-envelope}から包絡の全次数flat性が得られ、
特殊fiberの三次元は正則なので中央三次商も全次数flatとなる。
証明書は選んだ八方向、全非零二次補正、rank及び再構成後の中心性を記録する。
\end{proof}

\begin{remark}[一次の実行済み検証と残る条件]
三乗根の量子so3点では、$t$ を中央変数とする四生成元の斉次表示を使う。
基本語を $x^ay^bz^ct^d$ とした六つの二次書換えの三次曖昧語は
$zyx,tyx,tzx,tzy$ の四つである。
\texttt{verify\_cubic\_obstruction.py} は、次数1--3の20個の巡回potential方向を斉次化し、
各方向について四つの一次重なりを $\mathbb Q(q)$、$q^2+q+1=0$ 上で計算する。
80個全てが零となることを厳密検証した。
これはこの20方向で三次障害の一次項が零であることを示す。
平方零底上の全次数平坦性は上の停止性の証明と合わせて従う。
指定potential chartの高次包絡平坦性は定理\ref{thm:artin-potential-envelope}で別途証明した。
全intrinsic変形のこの表示による被覆は、
この有限計算からは結論しない。
命題\ref{prop:four-cubic-obstructions}と補題\ref{lem:flat-envelope-principal}自体は通常の証明であり、
有限標本で置き換えていない。
\end{remark}


\subsection{二次包絡の三次平坦性を支配する四係数}
\label{sec:quadratic-four-obstructions}
前節の再構成に残る条件のうち、三次の条件を有限行列として分離できる。
以下では全次数平坦性を三次平坦性から無条件に推論しない。

\begin{proposition}\label{prop:four-cubic-obstructions}
$A$ を剰余体 $k$ を持つ可換局所Artin環とし、$R_A$ を
定理\ref{thm:artin-section-ring}の平坦な切断環とする。
$V=R_{A,1}$、$I_n=\ker(V^{\otimes n}\to R_{A,n})$、
$B_A=T_A(V)/(I_2)$ と置く。
特殊fiberは $R_0=B_0/(G_0)$ で、$B_0$ は14641型Koszul AS4、
$G_0$ は正則三次元とする。
このとき、基底の選択後、四つの係数 $o_1,\ldots,o_4\in\mathfrak m_A$ があり、
\[
 (B_A)_3\simeq A^{19}\oplus A/(o_1,\ldots,o_4).
\]
従って $(B_A)_3$ が $A$-flatであることと、四係数が全て零であることは同値である。
イデアル $(o_1,\ldots,o_4)$ は基底によらず
\[
 \operatorname{Fitt}_0\bigl(I_3/(V I_2+I_2V)\bigr)
\]
である。特にこの障害イデアルの形成は、この有限表示の任意の底変更と両立する。
\end{proposition}
\begin{proof}
$V,I_2,I_3$ はそれぞれrank $4,6,45$ の有限自由加群である。
二次関係挿入写像
\[
 d:(V\otimes I_2)\oplus(I_2\otimes V)\longrightarrow I_3
\]
は $45\times48$ 行列で表示される。
特殊fiberで像の次元は $64-\dim(B_0)_3=44$ である。
従ってある $44\times44$ 小行列式は単元であり、可逆な行・列操作で
\[
 d\ \sim\
 \begin{pmatrix}I_{44}&0\\0&(o_1\ o_2\ o_3\ o_4)\end{pmatrix}
\]
とできる。剰余体上のrankが44なので $o_i\in\mathfrak m_A$ である。
これは $I_3/\operatorname{im}d\simeq A/(o_1,\ldots,o_4)$ とFittingイデアルの式を与える。
$R_{A,3}\simeq A^{19}$ なので
\[
 0\longrightarrow I_3/\operatorname{im}d\longrightarrow(B_A)_3
 \longrightarrow R_{A,3}\longrightarrow0
\]
は $A$-加群として分裂する。
最後に、局所Artin環上の有限flat加群は自由であり、
$A/J$ がflatならその剰余体上のrankは1なので $A/J\simeq A$、従って $J=0$ である。
逆向きは明らかである。
\end{proof}

\begin{lemma}[全次数平坦性の後の三次核]\label{lem:flat-envelope-principal}
上の設定で、さらに $B_A$ が次数別に $A$-flatであると仮定する。
$g\in\ker(B_A\to R_A)_3$ を $G_0$ の持上げとすると、
\[
 \ker(B_A\to R_A)=B_Ag=gB_A
\]
であり、$g$ は左右正則正規である。
従って全次数平坦性が得られた後には、三次核の正則正規主イデアル性を
別の無限個の方程式として課す必要はない。
\end{lemma}
\begin{proof}
$K_n=\ker((B_A)_n\to R_{A,n})$ と置く。
両辺が有限自由で右辺への写像は全射なので $K_n$ も有限自由で、底変更に適合する。
そのrankは $\dim(B_0)_n-\dim(R_0)_n=\dim(B_0)_{n-3}$ である。
左右の乗法写像
\[
 (B_A)_{n-3}\xrightarrow{\ b\mapsto bg\ } K_n,
 \qquad
 (B_A)_{n-3}\xrightarrow{\ b\mapsto gb\ } K_n
\]
は特殊fiberで $G_0$ の正則性と正規性によって同型である。
同rank有限自由加群間の写像なので、行列式が単元となり双方とも同型である。
全次数でこれを適用すれば主張を得る。
\end{proof}

\begin{proposition}[全20方向の同時一次包絡は全次数で平坦]
\label{prop:square-zero-potential-pbw}
$k$ は標数零で原始三乗根 $q$ を含むとし、
$A=k[e_1,\ldots,e_{20}]/(e_1,\ldots,e_{20})^2$ とする。
量子so3点のpotential $\Phi_0$ に、次数1--3の20個の巡回語をそれぞれ係数 $e_j$ で加える。
中央変数 $t$ で斉次化した六つの二次関係の包絡 $B_A$ は、
\[
 \{x^ay^bz^ct^d:a,b,c,d\ge0\}
\]
を $A$-基底に持つ。従って全次数で有限自由であり、
相対Koszul解消のrank列は $(1,4,6,4,1)$、双対cohomologyは最高次数の $A(4)$ のみである。
これはArtin底上の相対AS4性であり、非被約底環を含む $B_A$ 自身の有限大域次元を意味しない。

さらにこの包絡中に中央三次元 $g$ があり、特殊fiberが正則元 $G_0$ なら、
$g$ は左右正則で $B_A/(g)$ は次数別に有限自由である。
特に既存の一次potential・中央三次incidenceに属する全方向は、
包絡と切断環の平坦性に関しては実際の同時一次変形として実現される。
全intrinsic H3の変形の被覆は別問題である。
\end{proposition}
\begin{proof}
単なる四つのrank条件の外挿ではなく、平方零係数による書換えの停止性を使う。
六つの逆順対の書換えを、定数部分と平方零部分に分ける。
$e_j$ を既に含む項には定数部分だけを適用する。二回目の変形項は零になる。
斉次次数を固定し、まず平方零係数を含む語を係数を含まない語より小さく順序付ける。
同じ係数層では、$x,y,z$ の個数、次に $x<y<z<t$ の辞書式順序を使う。
定数部分の書換えはこの順序を狭義に下げ、変形項は係数層を下げる。
固定次数の語は有限個なので、全ての書換えは停止する。

左辺は六つの相異なる長さ2の逆順対であり、包含曖昧性はない。
重なる曖昧性は $zyx,tyx,tzx,tzy$ の四つだけである。
その定数部分は既知の量子so3 PBW恒等式、一次部分は
\texttt{verify\_cubic\_obstruction.py}の全80恒等式で零になる。
二次以上の係数積は底環で零である。
離れた二箇所の書換えは可換なので、全局所曖昧性が解消する。
停止する線型書換えのDiamondの議論により、既約語は主張した $A$-基底となる。

相対Koszul性も全次数flat性から確認できる。
augmentation $B_A\to A$ の被約bar複体の内部次数 $n$ 部分は有界有限自由 $A$-複体である。
特殊fiberのcohomologyは $B_0$ のKoszul性によりhomological次数 $n$ のみにあり、
そのrankは $n=0,1,2,3,4$ で $1,4,6,4,1$、それ以外で零である。
補題\ref{lem:artin-perfect}の議論を各内部次数に適用すると、相対Torも同じ対角に有限自由として集中する。
この基底を順に持ち上げて極小相対自由解消を構成すれば、
$B_A,B_A(-1)^4,B_A(-2)^6,B_A(-3)^4,B_A(-4)$ の長さ4の解消を得る。
その双対複体の各内部次数も有界有限自由である。
特殊fiberのAS Gorenstein性と同じ最小複体の議論から、双対cohomologyは最高次数の $A(4)$ のみとなる。

最後に各次数の $g$ による乗法は、有限自由加群間の写像で特殊fiberが単射である。
従って適切な最大小行列式が単元で、写像は分裂単射、余核は有限自由である。
左右両方に適用して正則性と商のflat性を得る。
\end{proof}

\begin{proposition}[巡回potential上の四syzygyは全Artin底で持ち上がる]
\label{prop:universal-four-syzygies}
量子so3点を特殊fiberとする任意の局所Artin底上の巡回potential変形について、
斉次二次包絡の三次部分は自由rank20である。平方零の仮定は不要である。
\end{proposition}
\begin{proof}
$\widetilde\Phi$ を四変数で三次に斉次化し、
$r_i=\partial_{x_i}\widetilde\Phi$、$r_t=\partial_t\widetilde\Phi$、
$c_i=[t,x_i]$ とする。六つの関係は $c_i,r_i$ で取れる。
$t$ を末尾へ移した表示との差は $c_i$ の線型結合なので、同じ二次包絡である。
二次式 $f=\sum a_{uv}uv$ に対し
\[
 D_t(f)=\sum a_{uv}(c_u v+u c_v),\qquad c_t=0
\]
と書く。自由代数で $D_t(f)=[t,f]$ なので、三つの恒等式
\[
 t r_i-r_i t-D_t(r_i)=0
\]
を得る。さらに巡回微分のEuler恒等式から
\[
 \sum_i(x_i r_i-r_i x_i)+D_t(r_t)=0
\]
を得る。これらは六関係の四つの線型syzygyであり、係数に打切りを課していない。

特殊fiberの64行48列の挿入行列はrank44で、上の四syzygyはその核の基底である。
後者は、$r_i$ の左係数に $t$ が現れる三本と、$x_i$ が現れる一本の独立性からも分かる。
\texttt{verify\_universal\_cubic\_syzygies.py}は全20係数の多項式環上で四恒等式を確認し、
特殊fiberの二行列のrank44と4も厳密確認する。
44単元minorによって挿入行列を $I_{44}$ と残り20行4列のSchur補行列へ変形する。
四syzygyの残り4成分は特殊fiberで基底なので可逆な4次行列をなし、
正確なsyzygy恒等式からこのSchur補行列は零である。
従って挿入行列の余核は $A^{20}$ である。
\end{proof}

\begin{theorem}[巡回potential chartの全Artin底での包絡平坦性]
\label{thm:artin-potential-envelope}
$A$ を剰余体 $k$ を持つ局所Artin $k$-代数とする。
$k$ は標数零で原始三乗根 $q$ を含むとし、
\[
 \Phi_A\in A\langle x,y,z\rangle_{\rm cyc}^{\le3},\qquad
 \Phi_A\bmod\mathfrak m_A=
 xyz-qxzy-\tfrac12(x^2+y^2+z^2)
\]
とする。三次に斉次化した $\widehat\Phi_A\in A[t]\langle x,y,z\rangle_{\rm cyc}$ に対し
\[
 B_A=A[t]\langle x,y,z\rangle/
 (\partial_x\widehat\Phi_A,\partial_y\widehat\Phi_A,
 \partial_z\widehat\Phi_A),\qquad \deg t=\deg x=\deg y=\deg z=1
\]
と置く。この二次包絡は全次数で有限自由、rankは $\binom{n+3}{3}$ であり、
形成は任意のArtin底変更と両立する。
$t$ は中央正則で、$B_A$ は左右noetherian、相対Koszul型14641・相対AS Gorensteinである。
中央三次元が与えられ特殊fiberで正則なら、その商も次数別有限自由となる。
全intrinsic変形がこのchartを通ること、全中央三次元の持上げ、H3族の大域的実現は本定理に含まない。
\end{theorem}
\begin{proof}
単に三次のflat性から高次を推論せず、potentialのdg複体を使う。
$A[t]$ 上の自由dg代数 $\mathcal D_A$ を、生成元
\[
 x_i:(\text{内部次数},\text{homological次数})=(1,0),\quad
 \xi_i:(2,1),\quad \zeta:(3,2)
\]
と微分
\[
 d x_i=d t=0,\qquad d\xi_i=\partial_{x_i}\widehat\Phi_A,
 \qquad d\zeta=\sum_i[x_i,\xi_i]
\]
で定める。$t$ は係数環の中央変数であるから、巡回微分恒等式により $d^2=0$ であり、
$H_0(\mathcal D_A)=B_A$ となる。
各内部次数 $n$ の部分は有界有限自由 $A$-複体である。

まず特殊fiber $\mathcal D_0$ を考える。
$t=0$ での商は斉次potential $xyz-qxzy$ の標準potential dg代数である。
このpotentialはCY3量子多項式環を定めるため、正のhomological次数でacyclicである。
ここで使う外部入力は \cite[\S7.1, (7.1.4)]{EG}
（そこで明示されたGinzburg Theorem5.3.1）である。
同箇所の定義・仮定を一次資料で確認しており、一般の非斉次potentialに逆向きを適用してはいない。
$\mathbb Q(q)$ 上の複体の正確性は複素数体への忠実平坦拡大から降下し、任意の $k$ へ拡大できる。

複体の短完全列
\[
 0\longrightarrow\mathcal D_0(-1)\xrightarrow{t}\mathcal D_0
 \longrightarrow\mathcal D_0/(t)\longrightarrow0
\]
の長完全列から、$j>0$ について $H_j(\mathcal D_0)=tH_j(\mathcal D_0)$ となる。
内部次数は下に有界なので、この等式は $H_j(\mathcal D_0)=0$ を意味する。
同時に $t$ は $H_0$ 上単射であり、$H_0/(t)$ は量子三変数多項式環となる。
従って $\dim H_0(\mathcal D_0)_n=\binom{n+3}{3}$ である。

任意の $n$ に対し、有限自由複体 $(\mathcal D_A)_n$ の特殊fiberのhomologyは次数零のみにある。
補題\ref{lem:artin-perfect}から、この複体は $A^{\binom{n+3}{3}}[0]$ と準同型である。
これが全次数flat性と底変更を同時に与える。底変更先が非flatでも、perfect複体から得た同定なので問題はない。
$t$ の正則性は次数ごとの有限自由加群間の写像を特殊fiberで判定して従う。
左右noetherian性は有限 $\mathfrak m_A$-進filtrationの商が $B_0$ の有限直和であることから従う。
相対Koszul解消と双対性は命題\ref{prop:square-zero-potential-pbw}のbar複体による議論を適用する。
その議論で平方零性を使ったのはflat性を得る最初の書換えだけであり、flat性を得た後は任意のArtin底で有効である。
最後の中央三次商も、同命題の分裂単射の議論で従う。
\end{proof}

\begin{proposition}[八方向の同時二次中心切断]
\label{prop:eight-second-order}
第\ref{sec:exact}節の $a=b=c=0$、切断係数 $(2,1,2,3,5)$ の点で、
\[
 A_2=k[h_0,\ldots,h_7]/(h_0,\ldots,h_7)^3
\]
上の巡回potentialと中央三次切断の同時変形が存在する。
その一次の八方向は座標変更と二つのスカラーを除いた一次incidenceの全体を張る。
包絡と三次商は全次数で $A_2$-flatである。
この二次jetの存在は、任意次数の同時形式積分や全intrinsic変形の被覆を意味しない。
\end{proposition}
\begin{proof}
\texttt{verify\_second\_order\_incidence.py}が有理係数を $\mathbb Q(q)$ 上で再生する。
既知七方向を14個の座標方向に順に加え、rankを増す七本を選び、
既存の横断方向を加える。全rankは22である。
この八方向について中央交換子の二次項を105行36列行列 $Q$ にまとめる。
元の一次行列を $M$ とすると
\[
 \operatorname{rank}M=\operatorname{rank}(M\mid-Q)=18
\]
である。従って $MH=-Q$ を厳密に解け、$H$ の40行36列が
potentialの20係数と三次元の20係数への二次補正を与える。

二次乗法では、逆順三対の係数行列も変化することに注意する。
それを $A_0+A_1+A_2$、既約語の係数を $C_0+C_1+C_2$ と書き、書換えを
\[
 B_0=-A_0^{-1}C_0,\quad
 B_1=-A_0^{-1}(C_1+A_1B_0),\quad
 B_2=-A_0^{-1}(C_2+A_1B_1+A_2B_0)
\]
で二次まで正しく展開する。
一次で正規化した関係をそのまま二次へ延長していない。
補正後の元関係からこの書換えを再構成し、三つの中央交換子の次数0,1,2を
全て再計算して零であることを確認する。
平方零係数のときと同じく、係数次数が上がるたび残りの許容次数が下がるため、この有限展開は停止する。
定理\ref{thm:artin-potential-envelope}から包絡の全次数flat性が得られ、
特殊fiberの三次元は正則なので中央三次商も全次数flatとなる。
証明書は選んだ八方向、全非零二次補正、rank及び再構成後の中心性を記録する。
\end{proof}

\begin{remark}[一次の実行済み検証と残る条件]
三乗根の量子so3点では、$t$ を中央変数とする四生成元の斉次表示を使う。
基本語を $x^ay^bz^ct^d$ とした六つの二次書換えの三次曖昧語は
$zyx,tyx,tzx,tzy$ の四つである。
\texttt{verify\_cubic\_obstruction.py} は、次数1--3の20個の巡回potential方向を斉次化し、
各方向について四つの一次重なりを $\mathbb Q(q)$、$q^2+q+1=0$ 上で計算する。
80個全てが零となることを厳密検証した。
これはこの20方向で三次障害の一次項が零であることを示す。
平方零底上の全次数平坦性は上の停止性の証明と合わせて従う。
指定potential chartの高次包絡平坦性は定理\ref{thm:artin-potential-envelope}で別途証明した。
全intrinsic変形のこの表示による被覆は、
この有限計算からは結論しない。
命題\ref{prop:four-cubic-obstructions}と補題\ref{lem:flat-envelope-principal}自体は通常の証明であり、
有限標本で置き換えていない。
\end{remark}

\section{三捩れ中心incidenceの八変数同時形式chart}
\label{sec:multivariate-formal-chart}
二次jetを更に延ばす代わりに、Rainsの中心式を多変数で正則化する。
この節で完成するのは指定potential・中央三次\emph{表示incidence}の形式近傍であり、
全intrinsic H3変形の被覆ではない。

\begin{lemma}[Rains式の印刷係数と独立検証]\label{lem:rains-corrected}
\cite[Appendix9.1, p.35]{EG}の中心式を、関係
\[
 XY-QYX-TZ^2+c_1Z+c_2,\quad
 YZ-QZY-TX^2+a_1X+a_2,\quad
 ZX-QXZ-TY^2+b_1Y+b_2
\]
で用いる際、$Y$ の一次係数中の
\[
 -T\bigl((Q^3b_2+2Q^2+QT^3+2Q+T^3+1)b_2
               +QT a_1c_1-T^2b_1^2\bigr)
\]
は、内側の余分な $b_2$ を除き、$Q^3b_2$ を $Q^3$ に置き換える。
この修正後の式を $\Psi$ とすると、全係数を変数とした
$\mathbb Q(Q,T)[a_1,b_1,c_1,a_2,b_2,c_2]$ 上で $\Psi$ は中心である。
この係数は著者サイトのMAGMA式
\url{https://math.mit.edu/~etingof/delpezzocenter}
及び本リポジトリの既存Clifford積分検証器の式と一致する。
本節ではその正しい式を別の三捩れ正規形で再検証したのであり、新しい中心式の発見とはしていない。
既存の $b_2=0$ の一変数検証には、この印刷係数の影響はない。
\end{lemma}
\begin{proof}
\texttt{verify\_rains\_central\_formula.py}に修正後の全式を文字通り収録した。
次数4以下の語に対し、元関係の一つの書換えから得る連立線型方程式を強連結成分ごとに解く。
既約語を $X^iY^jZ^k$ とし、係数体上で得た正規形が元関係の全102挿入を消すことを再検証する。
その上で三つの交換子を計算し、全て零となることを確認する。
全正規形の分母は $Q^2+Q+1=0,T=0$ で非零であることも確認する。
これは有限標本でなく全係数での恒等式である。
印刷式と修正式の差は $-TQ^3b_2(b_2-1)Y$ であり、同じ正規形で
その $X$ との交換子が非零となることも確認する。
従って印刷式をそのまま一般の $b_2$ で使用していない。
\end{proof}

\begin{theorem}[八変数の全次数平坦な中心切断形式族]
\label{thm:eight-formal-central-lift}
$q_0^2+q_0+1=0$ とし、$\eta_0^3=2,\ \zeta_0^3=3$ を選ぶ。
\[
 \Lambda=k[[\lambda,\kappa-2,\eta-\eta_0,\zeta-\zeta_0,a,b,c,D-5]]
\]
上でpotentialを
\[
 \Phi=(1+\kappa\lambda)xyz-q_0xzy+
 \lambda(x^3+\eta^3y^3+\zeta^3z^3)
 -\tfrac12(x^2+y^2+z^2)-ax-by-cz
\]
と置く。その斉次二次包絡には中央三次元 $G$ の同時形式持上げがあり、
$\lambda=0$ で
\[
 G=\kappa\Omega(a,b,c)+X+\eta^3Y+\zeta^3Z+D
\]
となる。ここで $X=x^3+q_0^2x$（$Y,Z$ も同様）、$\Omega(a,b,c)$ は既知のCasimirである。
全Artin商上で包絡と三次商は全次数有限自由であり、包絡は相対Koszul型14641である。
係数は以下の明示的な有理式の正則化で与えられ、任意次数への逐次的な存在仮定を必要としない。
\end{theorem}
\begin{proof}
$p=1+\kappa\lambda$ として補題\ref{lem:rains-corrected}へ
\[
 Q=\frac{q_0}{p},\qquad T=-\frac{3\lambda\eta\zeta}{p},
\]
\[
 (a_1,b_1,c_1)=\left(-\frac{\eta\zeta}{p},-\frac{\zeta}{\eta p},-\frac{\eta}{\zeta p}\right),
 \qquad
 (a_2,b_2,c_2)=\left(-\frac{\eta\zeta a}{p},-\frac{\zeta b}{p},-\frac{\eta c}{p}\right)
\]
を代入し、$X=x,Y=\eta y,Z=\zeta z$ と戻す。
この変換は表示potentialを非零定数倍するだけなので、同じ三関係を与える。
$\Psi$ の既約語正規形をこの代入後に $N$ と書く。
$N$ の定数係数を $N_\varnothing$、$x^3$ 係数を $N_{x^3}$ とする。
全ての非定数係数は $\lambda^2$ で割れ、さらに
\[
 \left.\frac{N_{x^3}}{\lambda^2}\right|_{\lambda=0}
   =9\eta^2\zeta^2(1+q_0)\ne0.
\]
従って
\[
 G=\frac{N-N_\varnothing}{N_{x^3}}+D
\]
は、共通因子 $\lambda^2$ を消去した後に $\Lambda$ 上正則である。

\texttt{verify\_central\_multivariate\_lift.py}はこの二重消滅と全ての先頭係数を厳密検証する。
具体的には、一般の正規形の各係数を $(Q-q_0,T)$ の次数2まで展開し、
次数0,1が\emph{各低次係数単項式ごとに}零であることを確認する。
従って分子の二重消滅は形式的な観察でなく、局所環のイデアル $(Q-q_0,T)^2$ への所属である。
分母の単元性は前補題で確認済みである。
二次の先頭式を代入して $N_{x^3}/\lambda^2$ で規格化すると、
全係数がちょうど $\kappa\Omega+X+\eta^3Y+\zeta^3Z$ となることも全変数で確認する。

包絡の各次数の形式完備化は定理\ref{thm:artin-potential-envelope}により有限自由である。
従って $\lambda$ は非零因子であり、$\lambda$ を反転して成立する中心恒等式は
$\lambda^2$ を消去した後にも成立する。
特殊fiberの $G$ は非零の中央三次元で、包絡は整域なので正則である。
同定理の中央三次商の議論を全Artin商へ適用して平坦性を得る。
\end{proof}

\begin{theorem}[指定表示incidenceの形式的滑らかさ]
\label{thm:formal-potential-incidence-smooth}
点 $\Phi_0,G_0=2\Omega+X+2Y+3Z+5$ における
20個のpotential係数と20個の三次係数の中央incidenceの完備局所環は、22変数形式冪級数環である。
特殊fiberで恒等となるaffine座標変更とpotential・切断のスカラーを除いた表示変形問題は、
定理\ref{thm:eight-formal-central-lift}の八変数形式族で表される。
有限の特殊fiber自己同型による大域群oidや、全intrinsic H3への形式versal性は含めない。
\end{theorem}
\begin{proof}
中央incidenceは次数4以下の有限個の係数方程式で定義される。
その一次行列は既検証の105行40列行列 $M$、rank18なので接空間次元は22である。
14個の座標・スカラー方向のrankは14である。
形式族の $a,b,c,\kappa,\eta,\zeta,D$ の七方向は既知七方向を張る。
特に $d(\eta^3),d(\zeta^3)$ は $\eta_0,\zeta_0\ne0$ により非零である。
残る $\lambda$ 方向のpotentialは
\[
 x^3+2y^3+3z^3+2xyz
\]
であり、既検証の横断方向に一致する。
中心三次の一次補正の選択差は固定potentialの中心方向なので、既知方向に吸収される。
従って八変数族と14変数の恒等近傍の座標作用から得る22変数形式写像は接空間上同型である。

この事実だけで形式的滑らかさを厳密に示せる。
incidenceの完備局所環を $R$ とすると、embedding dimension22より
$k[[z_1,\ldots,z_{22}]]\twoheadrightarrow R$ がある。
上の形式写像との合成
\[
 k[[z_1,\ldots,z_{22}]]\longrightarrow R\longrightarrow k[[s_1,\ldots,s_{22}]]
\]
は一次項が可逆なので形式逆写像定理により同型である。
従って最初の全射の核は零であり、$R$ は形式冪級数環となる。
この同型は座標の14方向と八変数sliceを同時に与える。
Lie安定化子は零なので、特殊fiberで恒等となる形式座標変換の商として八変数sliceを得る。
この最後の議論を、全intrinsic変形と表示変形の同一視に置き換えてはいない。
\end{proof}

\section{平坦なAbel変形のsmooth proper enhancement}
\label{sec:artin-enhancement}
前節の全次数平坦な切断環から、相対圏への比較を進める。
この節では逆Serre作用と次数シフトの同一視まで無条件に結論しない。

\begin{lemma}[nilpotent底上のcompact性の持上げ]\label{lem:dg-compact-lift}
$A$ を剰余体 $k$ の局所Artin環、$C$ を $A$-線型dg代数、$M\in D(C)$ とする。
$C_0=C\otimes_A^{\mathbf L}k$、$M_0=M\otimes_A^{\mathbf L}k$ と置く。
$M_0$ が $D(C_0)$ でcompactなら $M$ は $D(C)$ でcompactである。
\end{lemma}
\begin{proof}
$A$ の極大イデアルによる有限filtrationを、$A$-加群の有限個の三角へ分解する。
任意の $N\in D(C)$ はそれらとderived tensorすることにより、
$C_0$-加群をrestrictionした対象
$N\otimes_A^{\mathbf L}(\mathfrak m^j/\mathfrak m^{j+1})$ の有限拡張として得られる。
各商は $k$ の有限直和なので、この対象は
$N\otimes_A^{\mathbf L}k$ の有限直和のrestrictionである。
任意の $L\in D(C_0)$ に対しderived adjunctionは
\[
 \operatorname{RHom}_C(M,\operatorname{Res}L)
  \simeq\operatorname{RHom}_{C_0}(M_0,L)
\]
を与える。右辺は $M_0$ のcompact性により任意直和と可換である。
上の有限三角を通じてこの性質を $N$ へ移せば、$M$ のcompact性を得る。
\end{proof}

\begin{corollary}[smooth性のArtin持上げ]\label{cor:dg-smooth-lift}
$C_0$ が $k$ 上smoothであれば $C$ は $A$ 上smoothである。
さらに $C\otimes_A^{\mathbf L}k$ が有限次元総cohomologyを持てば、$C$ は $A$ 上properである。
\end{corollary}
\begin{proof}
smooth性は、前補題を $C^{\rm op}\otimes_A^{\mathbf L}C$ と対角加群 $C$ に適用する。
そのderived特殊fiberは $C_0$ の対角である。
proper性については、$A$ の有限filtrationから、$C$ のcohomologyは有界かつ有限生成となる。
局所Artin環上の有界上の有限自由分解から単元微分を除いた最小複体を取る。
特殊fiberが有界有限次元なので、その最小複体は有限個の有限自由項のみとなる。
従って $C\in\operatorname{Perf}(A)$ である。
\end{proof}

\begin{theorem}[平坦な切断環のqgrのenhancement]\label{thm:flat-qgr-enhancement}
$R_0$ を本稿のH3点の反標準環とし、$\mathcal G_0=\operatorname{QGr}R_0$ とする。
そのcompact圏がsmooth properなenhancementを持つと仮定する。
本稿の滑らかな中心・滑らかな楕円分岐を持つ点では、この既存のorder認識を適用する。
$A$ を局所Artin $k$-代数とし、$R_A$ を $R_0$ の次数別有限自由なconnected graded変形で、
一次生成・左右noetherianとする。
このとき $\mathcal G_A=\operatorname{QGr}R_A$ は $\mathcal G_0$ のflat Abel変形である。
$E_{i,A}=\pi R_A(i)$ は $E_{i,0}$ のflatな持上げであり、$D(\mathcal G_A)$ をcompactに生成する。
ある有限集合 $I\subset\mathbb Z$ について
\[
 E_A=\bigoplus_{i\in I}E_{i,A}
\]
はcompact generatorとなり、そのdg自己準同型代数は $A$ 上smooth properである。
従って $D(\mathcal G_A)^c$ はsmooth properな $A$-線型dg enhancementを持ち、
その特殊fiberは $\operatorname{Perf}(\mathcal G_0)$ である。
全ての $E_{i,A}$ 間のderived Homは $A$-perfectで、そのderived特殊fiberは元のderived Homである。
\end{theorem}
\begin{proof}
$R_A$ に付随する有向 $\mathbb Z$-代数は $R_0$ の付随代数のflat線型変形である。
$E_{i,0}$ の正方向Ext消滅は\eqref{eq:artin-forward-vanishing}で確認済みである。
\cite[Proposition4.4, Theorem4.5]{DeDekenLowen}により対応するflat Abel変形は
まさに $\operatorname{Qmod}$ であり、ここでは $\operatorname{QGr}R_A$ と一致する。
同定理の逆構成で $E_{i,A}$ は一意なflat持上げである。
Proposition4.6は局所有限性・有限生成性・noetherian性の持上げとも両立する。

特殊fiberの $E_{i,0}$ はcompact生成列である。
compact性はH3の有限cohomological次元と局所有限AS Cohen--Macaulay切断のcohomology計算、
生成性はAZ amplenessから従う。
\cite[Propositions5.8,5.9]{DeDekenLowen}をflat対象 $E_{i,A}$ に適用すれば、
compact性と生成性が持ち上がる。
この二命題の仮定は一般のflat Grothendieck圏変形であり、geometric helixの仮定を課していない。
同原論文pp.499--500の命題と証明を今回確認した。有限完全例外列を使用していない。

特殊fiberのsmooth proper圏には有限compact generatorがある。
compact生成列 $\{E_{i,0}\}$ の有限部分の厚い閉包にそのgeneratorが入るので、
ある有限直和 $E_0$ 自体がgeneratorとなる。
同じ命題5.9をこの有限直和に適用して $E_A$ を得る。

derived Homの底変更は、compactな第一引数に対する
\cite[Proposition5.5]{DeDekenLowen}とderived adjunctionにより
\[
 \operatorname{RHom}_{\mathcal G_A}(E_{i,A},E_{j,A})
       \otimes_A^{\mathbf L}k
 \simeq
 \operatorname{RHom}_{\mathcal G_0}(E_{i,0},E_{j,0})
\]
となる。右辺は有界有限次元である。
前系のperfect性の議論により左側のHom複体は $A$-perfectである。
標準のdg enhancementで $C_A=\operatorname{RHom}(E_A,E_A)$ を取ると、
$C_A\otimes_A^{\mathbf L}k\simeq C_0$ で、$C_0$ はsmooth properである。
系\ref{cor:dg-smooth-lift}から $C_A$ もsmooth properとなる。
compact generatorによるderived Morita同値で結論を得る。
\end{proof}

\begin{corollary}\label{cor:formal-family-enhanced}
定理\ref{thm:eight-formal-central-lift}の八変数族の各Artin商は、
上記の意味でsmooth proper enhancementを持つflat Abel変形を与える。
従って既存の全Artin単射性定理に必要なenhancement・perfect Hom・特殊fiber底変更の条件は、
この表示形式族では追加仮定でなく満たされる。
相対Serreと次数シフトの同一視は次の補題で確認する。全intrinsic H3変形の被覆と相対代数性は別問題である。
\end{corollary}

\begin{lemma}[次数シフトと逆Serre作用の剛性]\label{lem:serre-shift-rigidity}
定理\ref{thm:flat-qgr-enhancement}の設定で、さらに $\HH^1(\mathcal C_0)=0$ とする。
$T_A=(1)$ を次数シフトとし、$\tau_A=S_A^{-1}[2]$ とすると、
\[
 \tau_A\simeq T_A
\]
という自然同型がある。従って $\tau_A$ はheartを保つ。
\end{lemma}
\begin{proof}
smooth proper enhancementの相対Serre作用は、compact generatorの自己準同型代数 $C_A$ の
相対双対bimoduleで与えられ、derived底変更と両立する。
このbimoduleは左右perfectで、特殊fiberでは可逆である。
左右のtensor随伴のunit・counitの錐は特殊fiberで零なので、$A$ の有限filtrationによるderived Nakayamaにより零である。
従って相対Serre作用も自己同値である。
特殊fiberでは $\tau_0\simeq T_0$ なので、$F_A=\tau_A T_A^{-1}$ は恒等の持上げである。
Morita模型で $F_A$ は可逆なperfect $C_A$-bimoduleとなり、その特殊fiberは対角 $C_0$ である。
局所Artin底を小拡大に分ける。
各小拡大で既存の二つの持上げの差は
\[
 \Ext^1_{C_0^{\rm op}\otimes C_0}(C_0,C_0)\otimes\ker(A'\to A)
 =\HH^1(\mathcal C_0)\otimes\ker(A'\to A)=0
\]
に属する。対角bimodule自身が持上げなので、対象の持上げの通常の障害理論
\cite{LowenObject}により、帰納的に $F_A$ は恒等bimoduleと同型である。
ここでは新しい持上げの存在を $\HH^2$ の消滅から主張しておらず、既存の二つの持上げの一意性だけを使う。
$T_A$ はAbel圏上の正確な自己同値なので、結論を得る。
\end{proof}

\begin{theorem}[八変数族のArtin相対H3認識]\label{thm:formal-family-H3}
既存の滑らかな中心・楕円分岐order認識と定理\ref{thm:HH8}の適用点
$G_0=2\Omega+X+2Y+3Z+5$ を取る。
定理\ref{thm:eight-formal-central-lift}の各Artin商のpointed Abel圏は、
次のArtin相対的意味で強い内在的H3を与える：
noetherianなflat Abel変形、smooth proper enhancement、
heartを保つ逆Serre作用、相対exceptionalな構造対象、AZ ample反標準列、
三次Hilbert値を持つ有限自由な反標準環と相対14641二次包絡を持つ。
この反標準環は構成した中央三次商 $R_A$ とgraded同型である。
rank-one torsion-free条件は幾何学的fiberの条件として満たされる。
非正則なArtin底上で絶対的大域次元2や整域性を主張する意味ではない。
\end{theorem}
\begin{proof}
包絡・商の全次数flat性と相対AS4性は既証明であり、左右noetherian性も持つ。
定理\ref{thm:flat-qgr-enhancement}からflat Abel圏とsmooth proper enhancementを得る。
定理\ref{thm:HH8}の $\HH^1=0$ と前補題から $\tau_A\simeq(1)$ となり、heartを保つ。
構造対象のderived自己準同型は $A$-perfectで、特殊fiberは $k[0]$ なので
補題\ref{lem:artin-perfect}から $A[0]$ である。

正方向Homの自然写像
$(R_A)_n\to\Hom(\pi R_A,\pi R_A(n))$ は同型である。
実際、De Deken--Lowenの定理4.5の逆構成による正方向Homの線型変形は、
もとの付随 $\mathbb Z$-代数に戻る。自然写像も特殊fiberで恒等であり、有限自由加群間なので同型である。
$\tau_A\simeq(1)$ を反復すれば、この同型はorbit積とgraded積に両立する。
従って反標準環は $R_A$ 自身である。
AZ amplenessは正の有限生成noetherian付随代数のtails構成による。
二次包絡は $Q(R_A)=B_A$ である。三次元による商は二次関係を変えないからである。
Hilbert値とfiberのrank-one torsion-free性は既知のH3特殊fiber及び次数別flat性から従う。
この認識ではZhang twistによるH3不変性や正常三次のtwist輸送を仮定していない。
\end{proof}

\begin{lemma}[potentialと中央一次直線の形式的回復]\label{lem:potential-recovery}
この点のArtin変形で、中央一次部分は厳密に $At$ である。
またdehomogenizedな三関係空間から、次数3以下の巡回potentialは定数を除き単元倍まで一意に回復される。
従って、特殊fiberで恒等なpointed Abel同値は、本稿の二つのpotential・中央三次表示について
affine座標変更とpotential・切断のスカラーによる同値を誘導する。
\end{lemma}
\begin{proof}
\texttt{verify\_potential\_recovery.py}は中央一次条件の40行4列行列のrankが3で、核が $kt$ であることを確認する。
$t$ 自体は全底で中央なので、単元3小行列式で残りを消去すれば核は $At$ に限られる。
従ってgraded同型はこの直線を保存し、$t=1$ でaffineな三変数座標変更となる。

同検証器は20個の巡回語を微分し、三関係による次数2以下の商へ写す30行20列行列がrank19で、
核が $k\Phi_0$ であることを厳密確認する。
二つのpotentialが同じ関係空間を与える場合、小拡大の核を係数とする差の巡回微分は元関係空間に入る。
従って差は $\Phi_0$ のスカラー倍であり、potentialの単元倍を調整して消せる。
小拡大に関する帰納法で一意性を得る。

pointed Abel同値は内在的Serre作用と構造対象を保つので、反標準環のgraded同型を誘導する。
前定理でこの環は表示した $R_A$ と同型である。
その二次包絡と三次核直線は内在的に決まり、三次核の基底の変更は単元倍に限られる。
これと上の回復を合わせると最後の主張を得る。
\end{proof}

\begin{corollary}[任意のAbel変形に対する適合enhancement]\label{cor:any-abel-enhanced}
上の特殊fiberに対する任意のflat Abel変形は、
定理\ref{thm:artin-injective}で要求した適合enhancementを持つ。
従って同定理の同型類上の単射性は、この点の全flat pointed Abel変形に適用できる。
\end{corollary}
\begin{proof}
De Deken--Lowenの定理4.5は任意のAbel変形に対し $E_{i,0}$ の一意なflat持上げを与える。
定理\ref{thm:flat-qgr-enhancement}のcompact生成性・perfect Hom・smooth性の証明は、
この持上げ列とその正方向Homだけを使い、線型変形が一周期graded環から来ることを使わない。
従って同じ証明をそのまま適用できる。
この段階では、一周期性や二次包絡の認識を任意のAbel変形へ先取りしていない。
\end{proof}

\begin{theorem}[選んだH3点の全Artin形式分類]\label{thm:intrinsic-formal-eight}
滑らかな中心・滑らかな楕円分岐の既存order同定と定理\ref{thm:HH8}の下で、
点 $G_0=2\Omega+X+2Y+3Z+5$ のflat pointed Abel変形を、
特殊fiberで恒等となるpointed exact同値で割ったArtin同型類の変形問題は、
定理\ref{thm:eight-formal-central-lift}の八変数形式族で与えられる。
全てのArtin底でこの族による実現が存在し、そのパラメータは恒等近傍の座標gaugeを除き一意である。
特にこの点の全Artin変形は、前述の相対的意味で内在的H3として認識される。
これは一つの形式近傍の分類であり、相対モジュライstackの代数性、
非自明な特殊fiber自己同型を含む全群oid、成分の大域閉包・全境界の分類を含まない。
\end{theorem}
\begin{proof}
八変数族は前定理までに適合したsmooth proper enhancementを持つH3のflat Abel変形となった。
表示incidenceの恒等gauge商からAbel変形への写像は、補題\ref{lem:potential-recovery}により
全Artin底で同型類上単射である。
さらに系\ref{cor:any-abel-enhanced}と定理\ref{thm:artin-injective}により、
Abel変形から圏のMorita変形への写像も同型類上単射である。
従って八変数形式族から圏変形への接写像は単射である。

定理\ref{thm:HH8}により圏変形は無障害で、miniversal底は八変数形式円板である。
構造対象はexceptionalなので、その持上げを加えても接空間と障害は変わらない。
従って上の接写像は次元から同型となる。
八変数形式族をこのminiversal底へ比較する形式写像は、一次項が可逆なので形式同型である。
より正確には、miniversal写像の形式smooth性を用いてArtin商ごとに比較を持ち上げ、
逆極限で得た二つの八変数冪級数環間の写像に形式逆写像定理を適用する。
従って構成した族は圏変形に対してversalであり、全Artin同型類を実現する。
既に示した同型類上の単射性から、そのパラメータも一意である。
$\HH^1=0$ は恒等近傍の自己同値の一意性に、$\HH^0=k$ によるスカラー高次自己同型は
ここで扱う同型類の変形問題に影響しないことに対応する。

任意のflat pointed Abel変形を取る。
そのenhancementの圏変形を上の形式族で実現し、特殊fiberで恒等なMorita同値を選ぶ。
全Artin単射性により、元のAbel変形はその族のAbel変形とpointed exact同値である。
族の各Artin商のH3認識は定理\ref{thm:formal-family-H3}で証明済みなので、結論を得る。
大域代数化や形式近傍の外の対象へは、この議論を適用していない。
\end{proof}

\section{有限型の表示族と三捩れ因子を越える楕円境界}
\label{sec:algebraic-boundary}
形式的な圏の認識と、代数的な表示族の存在は区別する。
明示中心式のおかげで、後者については有限個の開条件へ進める。

\begin{proposition}[有限型底上の全次数平坦な表示]\label{prop:algebraic-presentation-flat}
定理\ref{thm:eight-formal-central-lift}の八変数有理係数表示は、
選んだ点を含む有限型smoothな八次元底のZariski開部分 $U$ 上に定義できる。
底はまず
\[
 \eta\zeta p(p^3+1)\ne0,\qquad p=1+\kappa\lambda
\]
に制限し、中心式の有限個の分母と、$N_{x^3}/\lambda^2$ をさらに可逆化する。
この $U$ 上で包絡と中央三次商は次数別有限局所自由で、任意の底変更に適合する。
各幾何学的fiberの包絡はnoetherian Koszul CY AS4の整域で、三次元は正則中央である。
この命題は全fiberのH3 admissibilityを結論しない。
\end{proposition}
\begin{proof}
中心式は有限個の有理係数を持ち、共通 $\lambda^2$ の消去は局所環の厳密な所属から得た。
従って形式冪級数の無限個の係数を代数化する必要はなく、有限個の分母を除く開部分で定義される。
残る $N_{x^3}/\lambda^2$ の $\lambda=0$ 値は
$9\eta^2\zeta^2(1+q_0)$ なので、必要な開部分は選んだ点を含む。

斉次先頭potentialを $X=x,Y=\eta y,Z=\zeta z$ で書くと、
三次元Sklyaninの三係数は
\[
 (p,-q_0,3\lambda\eta\zeta)
\]
となる。最初の二係数は非零で、その立方は $p^3$ と $-1$ で異なる。
\cite[Theorem1.1, Corollary6.1]{IS}により、全fiberでKoszul・polynomial Hilbert・CY3である。
同一次本文の仮定を今回確認した。
標準の三次元AS正則代数のnoetherian・domain定理も、従来と同じ既知入力として使う。
これらの性質をISのHilbert計算だけから推論してはいない。
\cite[Theorem3.3.2(i)]{EG}で任意の二次・一次potential項を加えたfiltered PBW変形を取り、
そのRees環と正則中央一次拡大の認識 \cite{LBS} により、fiberの包絡は主張したAS4となる。

全次数での\emph{底上の}flat性も説明する。
次数 $n$ の表示行列は有限行列で、全fiberの商次元は $\binom{n+3}{3}$ で一定である。
底 $U$ はreduced Jacobson環上の開部分なので、rankより一つ大きい小行列式は恒等的に零であり、
rankを与える小行列式は $U$ を覆う。
従って各次数の商は有限局所自由である。この議論に次数ごと異なる開部分を追加していない。

中心恒等式は稠密開上で成立し、包絡の各次数は $U$ 上torsion-freeなので $U$ 全体へ延長する。
既存の正規形の分母を除いたため、$x^3$ 係数1の三次元は各fiberで非零である。
fiberの包絡は整域なので正則である。
各次数の中央三次乗法は有限局所自由加群間でfiberごとに単射となり、
その余核も有限局所自由である。これが商のflat性と任意底変更を与える。
\end{proof}

\begin{proposition}[正規化した楕円曲線の明示式]\label{prop:algebraic-boundary-cubic}
\[
 C=\kappa+\kappa^2\lambda+
       \left(\frac{\kappa^3}{3}+9\eta^3\zeta^3\right)\lambda^2,
 \qquad
 F=q_0p(x^3+\eta^3y^3+\zeta^3z^3)+Cxyz
\]
と置く。
斉次先頭環のmultilinear関係行列
\[
 M=\begin{pmatrix}
 3\lambda x&-q_0z&py\\
 pz&3\lambda\eta^3y&-q_0x\\
 -q_0y&px&3\lambda\zeta^3z
 \end{pmatrix}
\]
は
\[
 \det M=3\lambda F
\]
を満たす。
従って $\lambda\ne0$ の点曲線の、$\lambda=0$ を越える正規化した三次族は $F=0$ である。
\[
 \Delta_E=C^3+27p^3\eta^3\zeta^3\ne0
\]
をさらに課せばこれはsmoothな相対楕円曲線となる。
選んだ点で $\Delta_E=170$ であり、この開条件は形式近傍を失わない。
$\lambda=0$ では関係行列は射影平面全体でrank2で、対応する写像は
\[
 (x:y:z)\longmapsto(x:q_0y:q_0^2z)
\]
である。選択した三次元の多重線型化は、$F|_{\lambda=0}$ の非零定数倍となる。
\end{proposition}
\begin{proof}
行列式と $3\lambda C=p^3-1+27\lambda^3\eta^3\zeta^3$ は直接展開である。
三つの立方項の係数が非零のHesse三次式について、三偏微分の積を使えば
特異点が存在する条件は $C^3+27p^3\eta^3\zeta^3=0$ と分かる。
$\lambda=0$ でこの式は $\kappa^3+27\eta^3\zeta^3$、指定点では $8+27\cdot2\cdot3=170$ となる。
\texttt{verify\_algebraic\_boundary.py}はこれらを全係数で検証する。
同検証器は $\lambda=0$ の各射影patchで全二次小行列式の零点がないこともGroebner基底 $(1)$ により確認する。
表示した対角写像を $M$ に代入すれば零となる。
三次元の多重線型評価では $xyz$ は $q_0^2xyz$、各三乗はそのままなので、
$\kappa q_0^2xyz+x^3+\eta^3y^3+\zeta^3z^3=q_0^{-1}F|_{\lambda=0}$ である。
この点曲線の明示化だけを、全相対THCR認識や全H3開性と同一視しない。
\end{proof}

\begin{theorem}[相対THCR境界と全fiberの整域性]\label{thm:relative-THCR-boundary}
前命題の $U$ を、$\Delta_E\ne0$ と、必要ならさらに選んだ点を含む有限の開部分へ制限する。
このとき $E=(F=0)\to U$ には関係行列から定まる相対自己同型 $\sigma$ があり、
$L=\OO_E(1)$ とすると
\[
 R_U/(t)\simeq B(E,L,\sigma)
\]
が任意の底変更に適合して成立する。
全幾何学的fiberで $t$ は $R$ 上正則中央であり、$R$ はnoetherian整域である。
この結論は全fiberのH3 smoothnessを含まない。
\end{theorem}
\begin{proof}
まず関係式を右側の点の座標に関してまとめると、もう一つの行列
\[
 M^{\rm r}=\begin{pmatrix}
 3\lambda x&pz&-q_0y\\
 -q_0z&3\lambda\eta^3y&px\\
 py&-q_0x&3\lambda\zeta^3z
 \end{pmatrix}
\]
を得て、その行列式も $3\lambda F$ である。
$\lambda=0$ では左右とも射影平面全体でrank2であり、左右の核は互いに逆の対角写像を与える。
$E\to U$ はprojectiveなので、左右のrank1以下の閉部分の像は閉である。
指定fiberにその像はないため、それらを除く有限の開部分は指定点を含む。

この開部分で $M$ の核は直線束となり、$\sigma:E\to\PP^2_U$ を与える。
$\lambda\ne0$ では $M^{\rm r}$ の行列式から像は $E$ に入り、
$E$ のflat性と $\lambda\ne0$ の稠密性からこの性質は全体へ延長する。
$M^{\rm r}$ から得る逆写像も同様で、核の一意性により互いに逆である。
$\lambda=0$ での式は既に検証した対角写像である。

この相対データからTHCRを作る。
各正次数のline bundle $L_n=L\otimes\sigma^*L\otimes\cdots\otimes(\sigma^{n-1})^*L$ はfiber次数 $3n$ であり、
$H^1$ が消えるので、その切断はrank $3n$ の有限局所自由加群で底変更に適合する。
THCRは次数1で生成される。fiber上での説明を与えると、$L$ のbase-point-freeな二切断のpencilを取り、
その完全列を他の正次数line bundleでtensorすれば、次数3より大きい場合の乗法全射を得る。
二つの次数3line bundleが非同型なら同じpencilの議論でよく、同型なら
平面三次曲線が二次関係を持たないことから乗法全射を得る。
従って生成性は全fiberで成立し、Nakayamaにより相対的にも成立する。
関係行列が $E$ と $\sigma$ のgraph上で消えるため、$S_U=B_U/(t)$ からTHCRへの自然全射がある。

中央三次元 $\overline G$ がこの全射で零になることを確認する。
まずgenericな $\sigma^3$ は恒等でない。
これは有理section $(0:\zeta:-\eta)$ の三反復を用いて厳密に検証できる。
\texttt{verify\_boundary\_graph.py}は
$\lambda=1/2,\kappa=2,\eta=\zeta=1$ でその軌道を計算し、
三反復後の第一座標が非零であることを確認する。
これはgenericな有理写像の恒等式を一つの正確な反例で否定しているのであり、全fiberを標本から外挿していない。

THCRは、generic曲線の関数体 $K(E)$ のskew Laurent環へ埋め込める。
次数1の切断の比は $K(E)$ を生成するため、非零な中央次数3元があれば $\sigma^3$ は関数体上恒等となる。
genericにはそうでないので、中央元 $\overline G$ のTHCRへの像は零である。
この像は $U$ 上の有限局所自由な三次切断加群のsectionなので、稠密generic点で零なら全体で零である。
従って全射は $S_U/(\overline G)$ を通る。
全fiberで $S$ はAS3整域、$\overline G\ne0$ なので、商のHilbert値は $1,3,6,9,\ldots$。
THCRのHilbert値と等しく、両辺は次数別有限局所自由なので、全射は各次数で同型である。

$t,G$ は各fiberの包絡上で正則中央列である。
実際 $t$ は包絡上正則であり、$S$ 上の $\overline G$ も正則である。
このため順序を入れ替えて $R$ 上の $t$ も正則となる。
商 $R/(t)$ はsmooth楕円曲線のTHCRなので整域であり、
connected gradingによる有限の $t$-進次数を取れば $R$ 自体も整域となる。
noetherian性は既に得たnoetherian包絡の商であることから従う。
\end{proof}

\begin{corollary}[幾何学的一般点のH3性]\label{cor:eight-generic-H3}
この明示八変数族の幾何学的一般点はH3である。
同時に、選んだ点での全Artin形式分類は定理\ref{thm:intrinsic-formal-eight}で与えられる。
これら二つの事実だけから、族の全H3 locusがZariski開であるとは結論しない。
\end{corollary}
\begin{proof}
$\lambda,\kappa,\eta,\zeta,a,b,c$ の関数体を $K$ とし、$D$ を独立な超越元とする。
非斉次包絡 $V_K$ はCY3 PBW環で有限大域次元を持ち、中心三次元 $H$ は $D$ に依存しない。
その最高記号は正次数の非零元なので、$H$ は $K$ 上超越的である。
\[
 V_K\otimes_KK(D)/(H+D)
 \simeq V_K\otimes_{K[H]}K(H)
\]
は中心局所化であり、有限大域次元を持つ。
前定理でsmoothな楕円THCR境界を認識したので、既存の補題\ref{lem:generic-D}と
入力\ref{ass:recognition}(3)を適用できる。
必要な分離的体拡大と幾何学的fiberの認識は同じ入力に従う。
これは独立な中心値 $D$ のgeneric点についての証明である。
任意の数値 $D$ や有限位数を横断する一様なH3開性を主張する議論ではない。
\end{proof}

\section{八次元の代数的表示stackとH3幾何学的矢}
\label{sec:presentation-stack}
この節では全H3 stackの代数性を仮定せず、構成した表示の有限型stackを先に作る。

\begin{lemma}[有限次数の平坦表示開部分]\label{lem:finite-presentation-open}
20個の次数1--3巡回potential係数の空間には、$\Phi_0$ を含むaffine座標変更とpotentialスカラーで不変なZariski開部分 $W$ があり、
その普遍斉次二次包絡の次数2,3,4部分はそれぞれrank10,20,35のvector bundleとなる。
\end{lemma}
\begin{proof}
各次数の自由語に対する関係挿入は、potential係数に多項式的な有限行列である。
定理\ref{thm:artin-potential-envelope}により、$\Phi_0$ の全Artin商でそのrankは期待値となる。
従ってrankより一つ大きい小行列式は、$\Phi_0$ の局所環の極大イデアルの全冪に入る。
Krull intersectionにより局所環で零である。
係数空間は整域なので、それらの多項式小行列式は恒等的に零となる。
各次数で期待rankの小行列式が少なくとも一つ非零となる条件は開で、$\Phi_0$ を含む。
三次数の共通開部分を取ればよい。
このrank条件は座標変更及びpotentialスカラーで不変である。
これは次数4以下だけの主張であり、ここで全potentialの全次数flat性を開性から先取りしていない。
\end{proof}

\begin{theorem}[smooth DMな八次元表示stack]\label{thm:presentation-DM-eight}
$\Phi_0,G_0=2\Omega+X+2Y+3Z+5$ を含む有限型smooth DM stack $\mathfrak P$ と、
明示八変数族の開部分からの\'{e}tale全射atlasが存在する。
$\mathfrak P$ の幾何学的対象は、指定したpotential chartの中央三次表示を
filtered affine座標変更と二つのスカラーで割ったものである。
全fiberの包絡はnoetherian Koszul CY AS4、商はnoetherian整域でsmooth楕円THCR境界を持つ。
$\mathfrak P$ のH3対象の間の全幾何学的pointed exact矢は、この表示の群作用で尽くされる。
ただし、H3対象が $\mathfrak P$ のZariski開部分をなすことや、全intrinsic H3族がここへ降下することは本定理に含まない。
\end{theorem}
\begin{proof}
前補題の $W$ 上の三次bundleを $V_3$、四次bundleを $V_4$ とする。
$V_3$ の全空間から零sectionを除き、三つの交換子
\[
 [G,x_i]=0\quad\text{in }V_4,\qquad i=1,2,3
\]
の零scheme $Z$ を取る。これは有限型schemeで、局所的に40変数105方程式のincidenceである。
作用群は
\[
 \Gamma=(\GL_3\ltimes\mathbb G_a^3)\times\mathbb G_{m,\Phi}\times\mathbb G_{m,G},
 \qquad\dim\Gamma=14
\]
である。affine代入で生じるpotentialの定数項は微分に寄与しないので捨てる。
この作用は普遍包絡及び交換子を保つ。

定理\ref{thm:formal-potential-incidence-smooth}から、$Z$ の指定点での完備局所環は22変数冪級数環である。
有限型の $k$-schemeなので、$Z$ はその点でsmoothな22次元である。
明示八変数有理族は有限開部分 $U$ 上で定義済みなので、作用による代数的写像
\[
 f:\Gamma\times U\longrightarrow Z
\]
がある。指定点で微分は同型であり、sourceとtargetはsmoothで同次元だから $f$ はそこで\'{e}taleである。
作用の同変性により、\'{e}tale locusは $\Gamma\times U'$ という形で、$U'$ は指定点を含む開部分である。
その像 $Z'=f(\Gamma\times U')$ は $\Gamma$-不変開部分となる。
必要な中心一次rank3及びpotential回復rank19の不変開条件も課す。
これらは指定点で検証済みなので、その点を失わない。

$\mathfrak P=[Z'/\Gamma]$ と置く。
$Z'$ は\'{e}tale全射によりsmooth22次元であり、$f$ の微分の単射性から作用のLie安定化子は零である。
従って商は有限型smooth DM stack、次元は $22-14=8$ である。
$U'\to\mathfrak P$ は $Z'$ への底変更後に $f$ となるので、\'{e}tale全射atlasである。
包絡・商・境界のfiberの性質は $U'$ 上で証明済みで、座標作用と幾何学的体拡大で保存されるから $Z'$ 全体へ移る。

最後にH3対象間の矢を確認する。
既存の中央楕円枝の認識により、pointed exact同値は反標準環とその二次包絡を復元する。
CY包絡の中央三次商では、この反標準環は表示した $R$ 自身である。
中央一次部分が $kt$ なので、graded同型はこの直線を保存し、$t=1$ でfiltered affine同型となる。
potential回復行列のrank19条件により、同じ二次関係空間を与えるpotentialはスカラー倍のみである。
三次核も一本なので、その基底の変更はもう一つのスカラーとなる。
従って全ての幾何学的矢は $\Gamma$ 作用から来る。
逆はgraded同型から明らかである。
次数全体の共通スカラーは $t=1$ の段階で消えており、余分な $\mathbb G_m$ 慣性を追加していない。
この幾何学的矢の結論を、未証明のH3 locusの開性や全圏族のfpqc降下に置き換えていない。
\end{proof}

\section{特異中心を持つH3の二つの因子}
\label{sec:nodal-centers}
本節では固定包絡 $B_{\mathbf0}$ の中心三次切断のH3判定
（$\Delta\ne0$ かつ超平面が中心三様体の射影双対に属さないこと）を、
既存の固定包絡分類定理
\texttt{papers/sections/cubic\_quantum\_triangle\_full\_cuts.tex} の
\texttt{thm:quantum-triangle-full-H3} からの入力として用いる。
この判定定理自体の全依存証明を今回新たに独立検証したわけではない。

$B_{\mathbf0}$ の非斉次中心は
\[
P=\Omega^3+f\Omega^2-XYZ+e(X^2+Y^2+Z^2)=0,
\quad f=\frac{1-q}{3},\quad e=-\frac{2q+1}{9}
\]
である。その三次斉次化が定める射影三様体 $\mathcal Y\subset\PP^4$ は
既存の分類により五つの通常二重点 $p_0,p_\eta$ を持つ。
$\Omega,X,Y,Z,t^3$ を座標とする $t=1$ のchartでは
\[
p_0=(0,0,0,0),\qquad
p_\eta=\left(-\frac{2f}{3},2e\eta_1,2e\eta_2,2e\eta_3\right),
\quad \eta_i\in\{\pm1\},\quad\eta_1\eta_2\eta_3=1.
\]
$\mathcal U\subset\PP^4$ はこの包絡の強いH3切断の許容開集合、
$L_i\subset\PP^4$ は切断超平面が $p_i$ を通ることを表す超平面とする。

\begin{theorem}[固定包絡H3の結節中心因子]
\label{thm:nodal-divisors}
上の既存H3判定定理の下で、$\mathcal U$ 内で中心三次曲面が特異になる
切断の集合は、正確に
\[
 \mathcal D_{\rm node}=\mathcal U\cap\bigcup_{i\in\{0\}\cup\{\eta\}}L_i
\]
である。五つの $L_i$ は線型一般位置にあり、
各 $\mathcal U\cap L_i$ は空でない既約な三次元因子で、
$\mathcal D_{\rm node}$ は $\mathcal U$ 上単純正規交叉因子となる。
一般点の中心はちょうど一つの $A_1$ 特異点を持ち、なお強いH3である。

$\mathfrak A_4$ は $L_0$ を固定し、残り四つの $L_\eta$ を推移的に置換する。
従って固定包絡の四次元DM stack $[\mathcal U/\mathfrak A_4]$ では
結節中心の因子にはちょうど二つの既約成分がある。
\end{theorem}
\begin{proof}
H3許容条件では切断超平面は $\mathcal Y$ の射影双対を避ける。
従って $\mathcal Y$ の滑らかな部分で接しない。
中心曲面の特異点は $\mathcal Y$ の五つの二重点の上に限られ、
そこで切断二次形式は非退化である。逆に $L_i$ 上の許容切断は
$p_i$ で通常二重点を持つ。よって集合としての等式を得る。

五点の射影座標を
\[
 (S,U,V,W,T)= (0,0,0,0,1),\quad
 (-1,\eta_1,\eta_2,\eta_3,1)
\]
とすれば、五行五列の座標行列の行列式は $16$ である。
従って五つの接続条件 $H(p_i)=0$ は独立な線型形式であり、
その零集合は座標超平面の配置と射影同値である。
各 $L_i\cap\mathcal U$ が非空であることは
命題\ref{prop:five-nodal-witnesses}で明示的に検証する。
$\mathcal U$ は開なので、交叉の局所横断性と既約性はそのまま保たれる。
最後に二符号変更と巡回置換からなる $\mathfrak A_4$ の作用は
四つの符号三つ組を推移的に置換し、原点を固定する。
\end{proof}

\begin{proposition}[五つの厳密な結節中心H3例]
\label{prop:five-nodal-witnesses}
\[
 G_D=2\Omega+X+2Y+3Z+D,
 \qquad
 \mathcal D_{\mathrm{ex}}=\{0\}\cup
 \left\{\frac{4f}{3}-2e(\eta_1+2\eta_2+3\eta_3):
 \eta_1\eta_2\eta_3=1,\ \eta_i=\pm1\right\}.
\]
$\mathcal D_{\mathrm{ex}}$ は相異なる五点からなり、各 $D\in\mathcal D_{\mathrm{ex}}$ について
$B_{\mathbf0}/(G_D)$ は強いH3である。
その中心曲面はただ一つの $A_1$ 特異点を持つ。
五つの構造対象付きAbel圏は互いにpointed exact同値でない。
\end{proposition}
\begin{proof}
$G_D=0$ から
$\Omega=-(X+2Y+3Z+D)/2$ を代入して得る三変数三次式を $F_D$ と置く。
$D=0$ では中心の二重点 $p_0$、その他の四値ではそれぞれ $p_\eta$ が
この切断上にある。$q^2+q+1=0$ を用いると、
各値の特異点Jacobiイデアルは正確に
\[
 (F_D,\partial_XF_D,\partial_YF_D,\partial_ZF_D)
   =(X-X_i,Y-Y_i,Z-Z_i) \subset\mathbb Q(q)[X,Y,Z]
\]
となる。右辺の点は $p_i$ の $X,Y,Z$ 座標である。
またその点の三変数Hessian行列式はどの場合も零でない。
これらは \texttt{computations/verify\_nodal\_cuts.py} が代数的数体上の
厳密Gr\"obner基底とHessianを計算し、
\texttt{nodal\_cuts\_certificate.json} に小証明書を出力して再現できる。
有限体への特殊化や浮動小数点判定は用いない。

全ての $D$ に対し境界の判別式は
$\Delta=2^3+27\cdot1\cdot2\cdot3=170\ne0$ である。
従って無限遠境界も滑らかであり、中心曲面の特異点は上記の
ただ一つの通常二重点で尽くされる。
その二次接錐は非退化なので、切断超平面は $\mathcal Y$ の
二重点を通る極限接超平面ではない。
滑らかな部分での接触も他の特異点を生じるため存在せず、
既存の固定包絡H3判定定理により強いH3となる。

五つの $D$ は互いに異なる。
固定包絡のpointed exact自己同型群 $\mathfrak A_4$ は
切断係数の $\alpha$ と $\varepsilon$ を保つ。
ここで $\alpha=2$ と規格化しているから、
異なる $D=\varepsilon$ の切断は同型にならない。
\end{proof}

\begin{remark}[変形と境界の解釈]
結節中心の五因子は許容H3開集合\emph{の内部}にある。
従って中心三次曲面の判別式をそのままH3成分の境界方程式とするのは誤りである。
一般の八次元表示stackでこの二つの因子がどのように延長・交叉し、
他の表示chartにどう現れるかは未解決である。
\end{remark}

\section{反標準環による包絡の一意性とAS曲面との比較}
\label{sec:as-surface-scope}

本節では、四次元AS正則環の三次商であることと、
三次元AS正則$\mathbb Z$-代数の$\qgr$であることを区別する。
後者がOkawa--Ueda \cite[Definition~1.1]{OUASSurfaces}のAS曲面であり、
del Pezzo型の標識には対応するvery strong helixが必要である。
以下の一意性は、構造対象を保つexact同値に対する主張である。

\begin{lemma}[飽和反標準環と二次包絡]\label{lem:anticanonical-quadratic-hull}
構造対象を持つ二つの圏$(\mathcal A,\OO)$、$(\mathcal A',\OO')$が
Serre作用を持ち、$\tau=S^{-1}[2]$、$\tau'=(S')^{-1}[2]$とする。
構造対象を保つexact同値は、反標準orbit環
\[
 R=\bigoplus_{n\geq0}\Hom(\OO,\tau^n\OO),\qquad
 R'=\bigoplus_{n\geq0}\Hom(\OO',(\tau')^n\OO')
\]
の次数付き同型を誘導する。
さらに$R=B/(G)$で、$B$が一次生成二次関係環、$G$が三次斉次元なら、
\[
 Q(R):=T_k(R_1)/\bigl(\ker(R_1\otimes_kR_1\longrightarrow R_2)\bigr)
 \simeq B.
\]
従って、このような二次包絡の次数付き同型類は$(\mathcal A,\OO)$から決まる。
\end{lemma}
\begin{proof}
exact同値の導来拡張はSerre作用をその一意性により保つ。
構造対象の同型と$F\tau\simeq\tau'F$の反復を用いると、
各次数のHomの同型がorbit積と両立する。
後半では、三次元による商は次数一・二の部分とその乗法を変えない。
よって$R_1=B_1$であり、二次乗法の核は$B$の二次関係空間そのものである。
$B$はこの関係空間で定義されるので結論を得る。
\end{proof}

この補題を斉次表示へ適用するには、表示した商環が
\emph{飽和反標準環そのものであること}を確認する必要がある。
任意の次数付き環の$\qgr$同値、任意の導来同値、
構造対象を指定しない同値からの包絡復元は、ここでは主張しない。

\begin{proposition}[AOUの無限位数blow-up表示との分離]
\label{prop:aou-infinite-order-exclusion}
式\eqref{eq:B}の$B_{\mathbf a}$と、そのH3切断
$R_{\mathbf a,u}=B_{\mathbf a}/(G_{\mathbf a,u})$を取る。
入力\ref{ass:recognition}(4)に従い、この商環が対応するpointed圏の
飽和反標準環であるとする。
このpointed圏は、Abdelgadir--Okawa--Ueda
\cite[\S9]{AOUCubic}で\emph{無限位数の平行移動}から構成される
反標準商$(\qgr F,\pi F)$とpointed exact同値ではない。
\end{proposition}
\begin{proof}
同文献の式(9.6)--(9.10)は$F=P/(C'_3)$という二次AS4包絡を与え、
Lemma~9.2は逆Serre作用を$(1)$と同定する。
AS Cohen--Macaulay性による飽和性と合わせて$Q(F)=P$を得る。
同節では中心一次直線は$kt$であり、$P/(t)$は滑らかな楕円曲線を
point schemeとする三次元Sklyanin環である。

一方、命題\ref{prop:pbw}より$Z(B_{\mathbf a})_1=kt$で、
\[
 B_{\mathbf a}/(t)=S_q
 =k\langle x,y,z\rangle/(yz-qzy,\ zx-qxz,\ xy-qyx).
\]
$q^3=1$かつ$q\ne1$なので、$S_q$の二次関係を多重線型化したschemeは
$\mathbb P^2$全体上の$\phi=\operatorname{diag}(1,q,q^2)$のgraphである。
実際、$(x',y',z')=(x,qy,q^2z)$は全関係を満たし、
任意の非零$(x,y,z)$でこの三線型方程式の係数行列はrank~2である。
したがってその核はこの点に限られ、標準的な三つのaffine chart上の
消去によりschemeとしてもgraphとなる。
特に$S_q$のpoint schemeは二次元の$\mathbb P^2$である。

もし上記のpointed exact同値があれば、補題
\ref{lem:anticanonical-quadratic-hull}により$B_{\mathbf a}\simeq P$となる。
唯一の中心一次直線は保存されるので$S_q\simeq P/(t)$を得るが、
次数付き同型はpoint schemeを保つため、上の次元差に矛盾する。
\end{proof}

\begin{remark}[除外の範囲]\label{rem:aou-exclusion-scope}
この命題は、AOUの\S9で指定された無限位数のblow-up表示を除外する。
有限位数の特殊化、別の構造対象や別の標識を持つAS曲面、
任意の導来同値を除外していない。
同文献Theorem~9.4はblow-up側とAS代数側の\emph{導来同値}であり、
ここで必要なpointed exact同値と置き換えることはできない。
また第\ref{sec:integrate}節の横断方向では中心一次商のpoint schemeも
楕円曲線へ変化するため、この障害を八次元変形族全体へ延長できない。
したがって、七次元族あるいは八次元族の一般元について、
「AS曲面である」「どの標識でもAS曲面にならない」のいずれも本節からは結論しない。
公表に向けた次の比較では、許容する構造対象・標識・同値の種類と、
AOU/Okawa--Ueda側の有限位数部分の適用範囲を先に固定する必要がある。
\end{remark}

\section{大域モジュライへの適用条件}\label{sec:component}
\begin{theorem}[相対モジュライ比較を明示した成分判定]\label{thm:conditional-component}
次の条件を仮定する。
\begin{enumerate}
\item 問題のH3対象を含む局所有限型のrigidified代数的モジュライstack $\mathfrak M$ があり、
上記の平坦な反標準環の族を受け入れる。
\item 一般点でH3の無限小変形から構造対象付き圏の無限小変形への写像が単射である。
\item 第\ref{sec:integrate}節の八次元一般表示族とその七次元族への特殊化が、
この相対モジュライ問題の族として実現される。
\end{enumerate}
このとき七次元族の一般点を通る既約成分は唯一で、次元は8である。
その点で $\mathfrak M$ は滑らかで、完成局所底は八次元形式円板となる。
\end{theorem}
\begin{proof}
条件(2)と定理\ref{thm:HH8}から接空間次元は8以下である。
条件(3)の一般表示族は、境界の平行移動を動かし、同型による冗長性を除いて八次元である。
その閉包が七次元族の一般点を含むので、局所次元は8以上となる。
一般慣性は自明であるため\'{e}tale局所環を取れ、接空間次元とKrull次元はともに8。
従って局所環は正則であり、唯一の既約成分を持つ。
\end{proof}

第\ref{sec:artin-reconstruction}節で条件(2)を適合したenhanced Abel変形の範囲では証明する。
他の条件を未検証のまま外して「全成分を完全に決定した」とはしない。
特に正則な八次元形式近傍を得ても、全既約成分の境界の記述はそれだけでは得られない。
完全決定には、一般族の像の閉包に含まれるH3対象を漏れなく認識し、
滑らかな楕円境界が消える点、中心が特異になる点、非PI領域、
別の表示を持つ点について全矢と慣性を求める必要がある。

有力な候補は、Etingof--Ginzburg型の一般中心三次商を含む八次元成分である。
Rainsの $E_6$ 型幾何学と比較する際には、
固定した楕円境界のfiltered deformationの六次元と、境界曲線・平行移動の二次元を区別する。
期待される重み付き射影表示やWeyl商を、そのまま全H3の同値と断定しない。
指定potential chartでの全Artin包絡認識と八変数の同時中心切断は後節で解決する。
全intrinsic形式変形との比較、H3族の相対実現、閉包の境界の認識は保持すべき課題である。
他方、指標多様体・特異点・実幾何との具体的接続は、それらの完成を待たずに調べられる。


\section{計算の役割と残る課題}
PBW還元、中心式、接空間行列、形式族の係数恒等式、結節例には厳密有理計算の検証器を付す。
再現入口は\path{computations/verify_all.py}であり、コードと小証明書を公開する。
計算が確認するのは記載した恒等式・階数・具体例である。
全次数の平坦性、圏の再構成、H3認識、代数的モジュライへの降下は、
本文の通常の証明と明示した入力に依存する。
投稿前に必要な核心は、外部H3認識の依存証明の再点検、
全intrinsic族の相対代数化・開性・降下、および既知の量子del Pezzo族との詳細な新規性比較である。
形式円板と表示stackの構成を、これら全ての解決として扱わない。
\UsePublicationReferences{AOUCubic,BG,CMR,CompanionBoundary,CompanionFricke,DeDekenLowen,EG,FIOS,IS,LBS,LVdB,LowenObject,NS,OUASSurfaces,Rains,Repo,StacksPerfect}
\begin{thebibliography}{99}
\small
\linespread{0.94}\selectfont
\setlength{\itemsep}{0pt}
\ifcsname pubref@DeDekenLowen\endcsname
\bibitem{DeDekenLowen} O. De Deken, W. Lowen,
\emph{Abelian and derived deformations in the presence of $\mathbb Z$-generating geometric helices},
J. Noncommut. Geom. 5 (2011), 477--505,
\href{https://doi.org/10.4171/JNCG/83}{doi:10.4171/JNCG/83}, Theorem 4.5, Proposition 4.6, pp. 495--496.
\fi
\ifcsname pubref@StacksPi1\endcsname
\bibitem{StacksPi1} The Stacks Project Authors,
\emph{The Stacks Project},
\href{https://stacks.math.columbia.edu/tag/0BTX}{Tag 0BTX}, Lemma 58.14.3
(arithmetic--geometric fundamental group exact sequence).
\fi
\ifcsname pubref@StacksPerfect\endcsname
\bibitem{StacksPerfect} The Stacks Project Authors,
\emph{The Stacks Project},
\href{https://stacks.math.columbia.edu/tag/0BDK}{Tag 0BDK}, Lemma 36.31.3.
\fi
\ifcsname pubref@EG\endcsname
\bibitem{EG} P. Etingof, V. Ginzburg,
\emph{Noncommutative del Pezzo surfaces and Calabi--Yau algebras},
J. Eur. Math. Soc. 12 (2010), 1371--1416.
\href{https://arxiv.org/abs/0709.3593v4}{arXiv:0709.3593v4}.
Theorem 3.3.2(i),(ii), Proposition 7.3.1, Appendix 9.1。
\fi
\ifcsname pubref@LBS\endcsname
\bibitem{LBS} L. Le Bruyn, S. P. Smith, M. Van den Bergh,
\emph{Central extensions of three dimensional Artin--Schelter regular algebras},
Math. Z. 222 (1996), 171--212, Theorem 2.6, Corollary 2.7。
\fi
\ifcsname pubref@BG\endcsname
\bibitem{BG} A. Braverman, D. Gaitsgory,
\emph{Poincar\'e--Birkhoff--Witt theorem for quadratic algebras of Koszul type},
J. Algebra 181 (1996), 315--328, Theorem 0.5。
\fi
\ifcsname pubref@IS\endcsname
\bibitem{IS} N. Iyudu, S. Shkarin,
\emph{Three dimensional Sklyanin algebras and Gr\"obner bases},
\href{https://arxiv.org/abs/1601.00564}{arXiv:1601.00564}, Theorem 1.1, Corollary 6.1。
\fi
\ifcsname pubref@NS\endcsname
\bibitem{NS} C. Negron, T. Schedler,
\emph{The Hochschild cohomology ring of a global quotient orbifold},
Adv. Math. 364 (2020), 106978.
\href{https://arxiv.org/abs/1809.08715}{arXiv:1809.08715}, Theorem 1.1, Appendix A。
\fi
\ifcsname pubref@FIOS\endcsname
\bibitem{FIOS} E. Faber, C. Ingalls, S. Okawa, M. Satriano,
\emph{On stacky surfaces and noncommutative surfaces},
\href{https://arxiv.org/abs/2206.13359v2}{arXiv:2206.13359v2}, 8 February 2024,
Theorem 1.5, Corollary 1.7。Brauer捩れ付きの比較として使用する。
\fi
\ifcsname pubref@FPSS\endcsname
\bibitem{FPSS} L. Fu, M. Porta, S. Scherotzke, N. Sibilla,
\emph{Hochschild--Kostant--Rosenberg isomorphism for derived Deligne--Mumford stacks},
\href{https://arxiv.org/abs/2509.00501v3}{arXiv:2509.00501v3}, 20 January 2026,
Proposition 6.17。一般root stackの恒等部分への帰着を比較し、$n=2$ に限る後半の式を $n=3$ へ流用しない。
\fi
\ifcsname pubref@TerwilligerDAHA\endcsname
\bibitem{TerwilligerDAHA} P. Terwilliger,
\emph{The universal Askey--Wilson algebra and DAHA of type $(C_1^\vee,C_1)$},
SIGMA 9 (2013), 047, \href{https://doi.org/10.3842/SIGMA.2013.047}{doi:10.3842/SIGMA.2013.047},
Definitions 2.1, 2.2, 2.6; Lemmas 2.3, 2.5, 2.8.
\fi
\ifcsname pubref@GoldmanTrace\endcsname
\bibitem{GoldmanTrace} W. M. Goldman,
\emph{Trace coordinates on Fricke spaces of some simple hyperbolic surfaces},
in \emph{Handbook of Teichm\"uller Theory}, Vol. II (2009), 611--684,
\href{https://doi.org/10.4171/055-1/16}{doi:10.4171/055-1/16},
\S5.2.1, equation (5.2.1.2).
\fi
\ifcsname pubref@GoldmanToledo\endcsname
\bibitem{GoldmanToledo} W. M. Goldman, D. Toledo,
\emph{Affine cubic surfaces and relative $\mathrm{SL}(2)$-character varieties of compact surfaces},
\href{https://arxiv.org/abs/1006.3838}{arXiv:1006.3838}.
Theorem 1: every smooth-at-infinity cubic with generic tritangent plane is a four-holed-sphere character variety.
\fi
\ifcsname pubref@SchockD4\endcsname
\bibitem{SchockD4} N. Schock,
\emph{The $W(E_6)$-invariant birational geometry of the moduli space of marked cubic surfaces},
Math. Nachr. 297 (2024), 2639--2666,
\href{https://doi.org/10.1002/mana.202300459}{doi:10.1002/mana.202300459}.
Section 2: tritangent triangle and its orthogonal $D_4$ root lattice.
\fi
\ifcsname pubref@IISDiscriminant\endcsname
\bibitem{IISDiscriminant} M. Inaba, K. Iwasaki, M.-H. Saito,
\emph{Dynamics of the sixth Painlev\'e equation},
S\'eminaires \& Congr\`es 14 (2006), 103--167,
\href{https://emis.de/ft/47232}{original text}, \S9.1, Definition 9.1.
Fricke判別式をトレースパラメータへ引き戻した場合の二乗因子について参照する。
本稿の$N_5^2M_5$との同一視は主張しない。
\fi
\ifcsname pubref@HavlicekPosta\endcsname
\bibitem{HavlicekPosta} M. Havl\'{\i}\v{c}ek, S. Po\v{s}ta,
\emph{Center of quantum algebra $U'_q(\mathrm{so}_3)$},
J. Math. Phys. 52 (2011), 043521,
\href{https://doi.org/10.1063/1.3579992}{doi:10.1063/1.3579992}.
三乗根を含むroot-of-unity中心の構造について参照し、規約の完全な同一視は追加検証事項。
\fi
\ifcsname pubref@CantatLorayCayley\endcsname
\bibitem{CantatLorayCayley} S. Cantat, F. Loray,
\emph{Dynamics on Character Varieties and Malgrange irreducibility of Painlev\'e VI equation},
Ann. Inst. Fourier 59 (2009), 2927--2978,
\href{https://doi.org/10.5802/aif.2512}{doi:10.5802/aif.2512}, \S1.5.
\fi
\ifcsname pubref@JangBoundary\endcsname
\bibitem{JangBoundary} S. u. Jang,
\emph{Tropical Dynamics of Markov Surfaces},
\href{https://arxiv.org/abs/2306.11357v2}{arXiv:2306.11357v2},
Lemma 4.8 and Theorem 4.10. 取得した2025年12月版の題名と番号を使用する。
\fi
\ifcsname pubref@MikoszWeberTriality\endcsname
\bibitem{MikoszWeberTriality} M. Mikosz, A. Weber,
\emph{Triality, characteristic classes, $D_4$ and $G_2$ singularities},
J. Homotopy Relat. Struct. 10 (2015), 995--1011,
\href{https://doi.org/10.1007/s40062-014-0094-3}{doi:10.1007/s40062-014-0094-3}, Proposition 3, \S5.
\fi
\ifcsname pubref@DixonElliptope\endcsname
\bibitem{DixonElliptope} K. Dixon,
\emph{The multi-moment maps of the nearly K\"ahler $S^3\times S^3$},
Geom. Dedicata 200 (2019), 351--362,
\href{https://doi.org/10.1007/s10711-018-0375-9}{doi:10.1007/s10711-018-0375-9}, Lemma 3.5.
\fi
\ifcsname pubref@HikamiCluster\endcsname
\bibitem{HikamiCluster} K. Hikami,
\emph{Note on character varieties and cluster algebras},
SIGMA 15 (2019), 003,
\href{https://doi.org/10.3842/SIGMA.2019.003}{doi:10.3842/SIGMA.2019.003}.
\fi
\ifcsname pubref@AWAvatars\endcsname
\bibitem{AWAvatars} N. Cramp\'e et al.,
\emph{The Askey--Wilson algebra and its avatars},
J. Phys. A 54 (2021), 063001,
\href{https://doi.org/10.1088/1751-8121/abd783}{doi:10.1088/1751-8121/abd783}.
\fi
\ifcsname pubref@CookeLacabanne\endcsname
\bibitem{CookeLacabanne} J. Cooke, N. Lacabanne,
\emph{Higher rank Askey--Wilson algebras as skein algebras},
Ann. Inst. Fourier 76 (2026), 1055--1117,
\href{https://doi.org/10.5802/aif.3729}{doi:10.5802/aif.3729}.
\fi
\ifcsname pubref@BousseauSkein\endcsname
\bibitem{BousseauSkein} P. Bousseau,
\emph{Strong positivity for the skein algebras of the $4$-punctured sphere and the $1$-punctured torus},
Comm. Math. Phys. 398 (2023), 1--58,
\href{https://doi.org/10.1007/s00220-022-04512-9}{doi:10.1007/s00220-022-04512-9}.
\fi
\ifcsname pubref@EOR\endcsname
\bibitem{EOR} P. Etingof, A. Oblomkov, E. Rains,
\emph{Generalized double affine Hecke algebras of rank 1 and quantized del Pezzo surfaces},
Adv. Math. 212 (2007), 749--796,
\href{https://doi.org/10.1016/j.aim.2006.11.008}{doi:10.1016/j.aim.2006.11.008}.
\fi
\ifcsname pubref@CMR\endcsname
\bibitem{CMR} L. Chekhov, M. Mazzocco, V. Rubtsov,
\emph{Quantised Painlev\'e monodromy manifolds, Sklyanin and Calabi--Yau algebras},
Adv. Math. 376 (2021), 107442,
\href{https://doi.org/10.1016/j.aim.2020.107442}{doi:10.1016/j.aim.2020.107442},
\href{https://arxiv.org/abs/1905.02772}{arXiv:1905.02772}.
原著v2のDefinition 1.1、式(1.5)、Theorem 1.5、Proposition 4.2を参照。
\fi
\ifcsname pubref@Rains\endcsname
\bibitem{Rains} E. M. Rains,
\emph{Filtered deformations of elliptic algebras},
\href{https://arxiv.org/abs/2107.13540v2}{arXiv:2107.13540v2}, 17 February 2022,
Propositions 2.2, 2.3, Example 2.1, Section 6末尾のConjecture 1とRemark。
\fi
\ifcsname pubref@LVdB\endcsname
\bibitem{LVdB} W. Lowen, M. Van den Bergh,
\emph{Hochschild cohomology of abelian categories and ringed spaces},
Adv. Math. 198 (2005), 172--221,
\href{https://arxiv.org/abs/math/0405227}{arXiv:math/0405227}。
\fi
\ifcsname pubref@LowenObject\endcsname
\bibitem{LowenObject} W. Lowen,
\emph{Hochschild cohomology, the characteristic morphism and derived deformations},
Compos. Math. 144 (2008), 1557--1580,
\href{https://arxiv.org/abs/0707.2602}{arXiv:0707.2602}。
\fi
\ifcsname pubref@Repo\endcsname
\bibitem{Repo} \emph{nc-cubic-moduli research repository},
\href{https://github.com/uedakazushi/nc-cubic-moduli}{GitHub main}。
本稿の依存する個別ファイル・ラベル・確認済み状態は研究ノートの入力表に固定する。
\fi
\ifcsname pubref@CompanionE6\endcsname
\bibitem{CompanionE6}\emph{三乗根量子三次曲面の正常形と形式変形},
日本語研究原稿、2026年11月18日、\PublicationVersion。
\href{https://github.com/uedakazushi/nc-cubic-e6/blob/main/papers/nc-cubic-e6.tex}{\texttt{nc-cubic-e6.tex}}。
\fi
\ifcsname pubref@CompanionFricke\endcsname
\bibitem{CompanionFricke}\emph{量子三次曲面のFricke幾何と周期},
日本語研究原稿、2026年11月18日、\PublicationVersion。
\href{https://github.com/uedakazushi/nc-cubic-e6/blob/main/papers/fricke-geometry.tex}{\texttt{fricke-geometry.tex}}。
\fi
\ifcsname pubref@CompanionBoundary\endcsname
\bibitem{CompanionBoundary}\emph{三次曲面のPell力学と境界特異点},
日本語研究原稿、2026年11月18日、\PublicationVersion。
\href{https://github.com/uedakazushi/nc-cubic-e6/blob/main/papers/cubic-boundary-prym.tex}{\texttt{cubic-boundary-prym.tex}}。
\fi
\ifcsname pubref@AOUCubic\endcsname
\bibitem{AOUCubic}
T. Abdelgadir, S. Okawa, K. Ueda,
\emph{Compact moduli of marked noncommutative cubic surfaces},
\href{https://arxiv.org/abs/2404.00175v1}{arXiv:2404.00175v1} (2024).
\S9, equations (9.1)--(9.10), Lemma 9.2, Theorem 9.4.
\fi
\ifcsname pubref@OUASSurfaces\endcsname
\bibitem{OUASSurfaces}
S. Okawa, K. Ueda,
\emph{Marked Artin--Schelter surfaces of del Pezzo types},
\href{https://arxiv.org/abs/2508.04986v1}{arXiv:2508.04986v1} (2025).
Definition 1.1.
\fi
\ifcsname pubref@BoundaryDHKK\endcsname
\bibitem{BoundaryDHKK} G. Dimitrov, F. Haiden, L. Katzarkov, M. Kontsevich,
\emph{Dynamical systems and categories},
\href{https://arxiv.org/abs/1307.8418}{arXiv:1307.8418},
Definitions 2.1, 2.4, Proposition 2.2, Lemma 2.3, Theorem 2.6.
\fi
\ifcsname pubref@BoundaryGoldmanNeumann\endcsname
\bibitem{BoundaryGoldmanNeumann} W. M. Goldman, W. D. Neumann,
\emph{Homological action of the modular group on some cubic moduli spaces},
Math. Res. Lett. 12 (2005), 575--591,
\href{https://doi.org/10.4310/MRL.2005.v12.n4.a11}{doi:10.4310/MRL.2005.v12.n4.a11},
\href{https://arxiv.org/abs/math/0402039}{arXiv:math/0402039},
Theorems 1--2, Proposition 3.
\fi
\ifcsname pubref@BoundaryDyckerhoff\endcsname
\bibitem{BoundaryDyckerhoff} T. Dyckerhoff,
\emph{Compact generators in categories of matrix factorizations},
Duke Math. J. 159 (2011), 223--274,
\href{https://doi.org/10.1215/00127094-1415869}{doi:10.1215/00127094-1415869},
\href{https://arxiv.org/abs/0904.4713v5}{arXiv:0904.4713v5},
Theorem 5.9, Corollary 6.5.
\fi
\ifcsname pubref@BoundaryHuaKeller\endcsname
\bibitem{BoundaryHuaKeller} Z. Hua, B. Keller,
\emph{Cluster categories and rational curves},
Geom. Topol. 28 (2024), 2569--2634,
\href{https://doi.org/10.2140/gt.2024.28.2569}{doi:10.2140/gt.2024.28.2569},
\href{https://arxiv.org/abs/1810.00749v6}{arXiv:1810.00749v6},
Theorem 5.9, pp.36--37.
\fi
\ifcsname pubref@BoundaryMatherYau\endcsname
\bibitem{BoundaryMatherYau} J. N. Mather, S. S.-T. Yau,
\emph{Classification of isolated hypersurface singularities by their moduli algebras},
Invent. Math. 69 (1982), 243--251,
\href{https://doi.org/10.1007/BF01399504}{doi:10.1007/BF01399504}.
\fi
\ifcsname pubref@BoundaryMumford\endcsname
\bibitem{BoundaryMumford} D. Mumford,
\emph{Prym Varieties I}, in \emph{Contributions to Analysis},
Academic Press (1974), 325--350,
\href{https://doi.org/10.1016/B978-0-12-044850-0.50032-0}{doi:10.1016/B978-0-12-044850-0.50032-0},
\S\S1--3.
\fi
\ifcsname pubref@BoundaryMirandaPersson\endcsname
\bibitem{BoundaryMirandaPersson} R. Miranda, U. Persson,
\emph{On extremal rational elliptic surfaces},
Math. Z. 193 (1986), 537--558,
\href{https://doi.org/10.1007/BF01160474}{doi:10.1007/BF01160474},
Table 5.2.
\fi
\ifcsname pubref@BoundaryHowald\endcsname
\bibitem{BoundaryHowald} J. Howald,
\emph{Multiplier Ideals of Sufficiently General Polynomials},
\href{https://arxiv.org/abs/math/0303203}{arXiv:math/0303203},
Theorems 11--12, Corollary 13.
\fi
\ifcsname pubref@BoundaryCasals\endcsname
\bibitem{BoundaryCasals} R. Casals,
\emph{Lagrangian skeleta and plane curve singularities},
J. Fixed Point Theory Appl. 24 (2022), 34,
\href{https://doi.org/10.1007/s11784-022-00939-8}{doi:10.1007/s11784-022-00939-8},
Theorem 1.1, Corollary 1.2, \S2.1.
\fi
\ifcsname pubref@BoundaryCamara\endcsname
\bibitem{BoundaryCamara} L. M. C\^amara,
\emph{On the classification of quasi-homogeneous curves},
\href{https://arxiv.org/abs/1009.1664v1}{arXiv:1009.1664v1},
Theorem 1(iii), \S4.3.
\fi
\ifcsname pubref@FrickeFriedman\endcsname
\bibitem{FrickeFriedman}
R. Friedman, \emph{On the geometry of anticanonical pairs},
\href{https://arxiv.org/abs/1502.02560v2}{arXiv:1502.02560v2} (2016),
Proposition 3.12.
\fi
\ifcsname pubref@FrickeBoalchPaluba\endcsname
\bibitem{FrickeBoalchPaluba}
P. Boalch, R. Paluba, \emph{Symmetric cubic surfaces and $G_2$ character varieties},
J. Algebraic Geom. 25 (2016), 607--631,
\href{https://doi.org/10.1090/jag/668}{doi:10.1090/jag/668},
Lemma 6, Corollary 11, Proposition 18, Remark 19.
\fi
\ifcsname pubref@PymR13\endcsname
\bibitem{PymR13} B. Pym,
\emph{Quantum deformations of projective three-space},
Adv. Math. 281 (2015), 1216--1241,
\href{https://doi.org/10.1016/j.aim.2015.06.005}{doi:10.1016/j.aim.2015.06.005},
\href{https://arxiv.org/abs/1403.6444}{arXiv:1403.6444}, \S3.4, (9).
\fi
\ifcsname pubref@CompanionAmbient\endcsname
\bibitem{CompanionAmbient}
\emph{非可換三次曲面のambient射と$R(1,3)$成分},
\href{https://github.com/uedakazushi/nc-cubic-ambient/blob/main/papers/nc-cubic-ambient.tex}
{nc-cubic-ambient research manuscript},
\texttt{thm:r13-affine}, \texttt{prop:center}, \texttt{sec:surface-fiber}.
引用範囲は次数付き中心拡大の同型商と中心三次ペンシルであり、
AOUからの支配性の追加仮定は用いない。
\fi
\end{thebibliography}

\end{document}
